\documentclass[11pt,letterpaper]{article}

\usepackage[margin=1in]{geometry}
\usepackage{amsmath,amssymb}
\usepackage{booktabs}
\usepackage{array}
\usepackage{graphicx}
\usepackage{float}
\usepackage{xcolor}
\usepackage{microtype}
\usepackage[numbers,sort&compress]{natbib}
\usepackage[font=small,labelfont=bf,skip=5pt]{caption}
\PassOptionsToPackage{hyphens}{url}
\usepackage[colorlinks=true,linkcolor=blue!50!black,citecolor=blue!50!black,urlcolor=blue!50!black]{hyperref}
\Urlmuskip=0mu plus 1mu\relax
\emergencystretch=1.3em
\renewcommand{\bibfont}{\footnotesize\raggedright}
\setlength{\bibsep}{1pt}
\usepackage[compact]{titlesec}
\setlength{\textfloatsep}{9pt plus 2pt minus 3pt}
\renewcommand{\topfraction}{0.85}
\renewcommand{\bottomfraction}{0.6}
\renewcommand{\textfraction}{0.1}
\renewcommand{\floatpagefraction}{0.8}

\newcolumntype{L}[1]{>{\raggedright\arraybackslash}p{#1}}

\newif\ifdarkmode
\ifdefined\DARKMODE \darkmodetrue \else \darkmodefalse \fi
\ifdarkmode
  \input{darkstyle}
\fi

\newcommand{\singleauthorinfo}[1]{\\[0.25em]{\small #1}}
\newcommand{\q}[1]{\texttt{#1}}
\newcommand{\fmt}[1]{\texttt{#1}}
\newcommand{\fig}[1]{\ifdarkmode figures/dark/#1\else figures/#1\fi}
\newcommand{\slot}[1]{$\langle$\textit{#1}$\rangle$}
% Generated by paper8/scripts/build_claims.py from testsuite/results/turing/longctx/ (and P6's generated
% claims for the section on Paper 6). Do not edit by hand.
\newcommand{\PeAcceptEE}{60.5}
\newcommand{\PeAcceptFF}{60.3}
\newcommand{\PeAcceptFp}{60.8}
\newcommand{\PeAdaBoostMemMhz}{9,001}
\newcommand{\PeAdaBoostOneFasterHi}{27}
\newcommand{\PeAdaBoostOneFasterLo}{19}
\newcommand{\PeAdaMemMhz}{7,001}
\newcommand{\PeAdaOneFasterHi}{16}
\newcommand{\PeAdaOneFasterLo}{10}
\newcommand{\PeAnswerMax}{107}
\newcommand{\PeAnswerMaxCalibTwo}{404}
\newcommand{\PeAnswerMaxNinetySix}{107}
\newcommand{\PeArmEESixtyFourAcc}{70.8}
\newcommand{\PeArmEESixtyFourDecode}{63.5}
\newcommand{\PeArmEESixtyFourHi}{79.0}
\newcommand{\PeArmEESixtyFourItemS}{115.1}
\newcommand{\PeArmEESixtyFourK}{68}
\newcommand{\PeArmEESixtyFourLo}{61.1}
\newcommand{\PeArmEESixtyFourN}{96}
\newcommand{\PeArmEESixtyFourPeak}{10,214}
\newcommand{\PeArmEESixtyFourPrefill}{573}
\newcommand{\PeArmEEThirtyTwoAcc}{77.1}
\newcommand{\PeArmEEThirtyTwoDecode}{71.0}
\newcommand{\PeArmEEThirtyTwoHi}{84.4}
\newcommand{\PeArmEEThirtyTwoItemS}{52.4}
\newcommand{\PeArmEEThirtyTwoK}{74}
\newcommand{\PeArmEEThirtyTwoLo}{67.7}
\newcommand{\PeArmEEThirtyTwoN}{96}
\newcommand{\PeArmEEThirtyTwoPeak}{8,971}
\newcommand{\PeArmEEThirtyTwoPrefill}{632}
\newcommand{\PeArmEFNinetySixAcc}{58.3}
\newcommand{\PeArmEFNinetySixDecode}{26.9}
\newcommand{\PeArmEFNinetySixFree}{615}
\newcommand{\PeArmEFNinetySixHi}{71.2}
\newcommand{\PeArmEFNinetySixItemS}{190.5}
\newcommand{\PeArmEFNinetySixK}{28}
\newcommand{\PeArmEFNinetySixLo}{44.3}
\newcommand{\PeArmEFNinetySixN}{48}
\newcommand{\PeArmEFNinetySixPeak}{10,649}
\newcommand{\PeArmEFNinetySixPrefill}{523}
\newcommand{\PeArmEFSixtyFourAcc}{82.3}
\newcommand{\PeArmEFSixtyFourDecode}{33.6}
\newcommand{\PeArmEFSixtyFourHi}{88.6}
\newcommand{\PeArmEFSixtyFourItemS}{116.0}
\newcommand{\PeArmEFSixtyFourK}{79}
\newcommand{\PeArmEFSixtyFourLo}{73.5}
\newcommand{\PeArmEFSixtyFourN}{96}
\newcommand{\PeArmEFSixtyFourPeak}{9,689}
\newcommand{\PeArmEFSixtyFourPrefill}{575}
\newcommand{\PeArmFFSixtyFourAcc}{67.7}
\newcommand{\PeArmFFSixtyFourDecode}{35.3}
\newcommand{\PeArmFFSixtyFourHi}{76.2}
\newcommand{\PeArmFFSixtyFourItemS}{116.2}
\newcommand{\PeArmFFSixtyFourK}{65}
\newcommand{\PeArmFFSixtyFourLo}{57.8}
\newcommand{\PeArmFFSixtyFourN}{96}
\newcommand{\PeArmFFSixtyFourPeak}{9,150}
\newcommand{\PeArmFFSixtyFourPrefill}{574}
\newcommand{\PeArmFFThirtyTwoAcc}{74.0}
\newcommand{\PeArmFFThirtyTwoDecode}{48.3}
\newcommand{\PeArmFFThirtyTwoHi}{81.7}
\newcommand{\PeArmFFThirtyTwoItemS}{54.5}
\newcommand{\PeArmFFThirtyTwoK}{71}
\newcommand{\PeArmFFThirtyTwoLo}{64.4}
\newcommand{\PeArmFFThirtyTwoN}{96}
\newcommand{\PeArmFFThirtyTwoPeak}{8,449}
\newcommand{\PeArmFFThirtyTwoPrefill}{616}
\newcommand{\PeArmFpThirtyTwoAcc}{77.1}
\newcommand{\PeArmFpThirtyTwoDecode}{9.3}
\newcommand{\PeArmFpThirtyTwoHi}{84.4}
\newcommand{\PeArmFpThirtyTwoItemS}{61.9}
\newcommand{\PeArmFpThirtyTwoK}{74}
\newcommand{\PeArmFpThirtyTwoLo}{67.7}
\newcommand{\PeArmFpThirtyTwoN}{96}
\newcommand{\PeArmFpThirtyTwoPeak}{10,225}
\newcommand{\PeArmFpThirtyTwoPrefill}{635}
\newcommand{\PeArmHours}{3.12}
\newcommand{\PeArmPeakMax}{10,649}
\newcommand{\PeArmTritonAcc}{70.8}
\newcommand{\PeArmTritonDecode}{34.8}
\newcommand{\PeArmTritonHi}{79.0}
\newcommand{\PeArmTritonItemS}{116.7}
\newcommand{\PeArmTritonK}{68}
\newcommand{\PeArmTritonLo}{61.1}
\newcommand{\PeArmTritonN}{96}
\newcommand{\PeArmTritonPeak}{10,231}
\newcommand{\PeArmTritonPrefill}{571}
\newcommand{\PeAuditChi}{7.0}
\newcommand{\PeAuditCloseN}{4}
\newcommand{\PeAuditCloseWord}{four}
\newcommand{\PeAuditFloorDrafted}{5.5}
\newcommand{\PeAuditFloorUndrafted}{3.6}
\newcommand{\PeAuditN}{11}
\newcommand{\PeAuditSessions}{13}
\newcommand{\PeAuditSlowHi}{13.6}
\newcommand{\PeAuditSlowLo}{9.2}
\newcommand{\PeBaselineMiB}{379}
\newcommand{\PeBenchBest}{469.8}
\newcommand{\PeBenchConfigs}{72}
\newcommand{\PeBenchDefault}{703.9}
\newcommand{\PeBenchFd}{145.2}
\newcommand{\PeBenchFdOverFloor}{1.3}
\newcommand{\PeBenchFloor}{111}
\newcommand{\PeBytesEE}{2,176}
\newcommand{\PeBytesEF}{1,664}
\newcommand{\PeBytesErrMax}{47}
\newcommand{\PeBytesErrMin}{42}
\newcommand{\PeBytesFF}{1,152}
\newcommand{\PeBytesFp}{4,096}
\newcommand{\PeCacheMiBEE}{36.1}
\newcommand{\PeCacheMiBEF}{27.6}
\newcommand{\PeCacheMiBFF}{19.1}
\newcommand{\PeCacheMiBFp}{68.0}
\newcommand{\PeCalMkFortyEight}{8/8}
\newcommand{\PeCalMkOneTwentyEight}{8/8}
\newcommand{\PeCalMkSixteen}{8/8}
\newcommand{\PeCalMvAcc}{62.5}
\newcommand{\PeCalMvFortyEight}{0/8}
\newcommand{\PeCalMvTwelve}{5/8}
\newcommand{\PeCalMvTwentyFour}{1/8}
\newcommand{\PeCalVtFour}{1/8}
\newcommand{\PeCalVtThree}{0/8}
\newcommand{\PeCalVtTwo}{1/8}
\newcommand{\PeCardMiB}{11,264}
\newcommand{\PeCensusBusy}{96}
\newcommand{\PeCensusFinMs}{21.4}
\newcommand{\PeCensusGemm}{15.8}
\newcommand{\PeCensusGemv}{53.2}
\newcommand{\PeCensusIdleHi}{1.8}
\newcommand{\PeCensusIdleLo}{1.5}
\newcommand{\PeCensusKernels}{869}
\newcommand{\PeCensusMatrix}{92}
\newcommand{\PeCensusMgemm}{23.0}
\newcommand{\PeCensusMs}{45.6}
\newcommand{\PeCfBuunAcc}{97.0}
\newcommand{\PeCfBuunDiff}{\ensuremath{-2.0}}
\newcommand{\PeCfBuunMtpZero}{40.91}
\newcommand{\PeCfBuunNdZero}{26.72}
\newcommand{\PeCfBuunOtt}{2}
\newcommand{\PeCfBuunP}{0.5}
\newcommand{\PeCfBuunRefAcc}{99.0}
\newcommand{\PeCfBuunTen}{0}
\newcommand{\PeCfFitWMtpDeep}{65.44}
\newcommand{\PeCfFitWMtpZero}{74.06}
\newcommand{\PeCfFitWNdDeep}{44.42}
\newcommand{\PeCfFitWNdZero}{48.53}
\newcommand{\PeCfFtwoMax}{+0.8}
\newcommand{\PeCfFtwoMin}{+0.3}
\newcommand{\PeCfGbsNif}{314}
\newcommand{\PeCfGbsW}{317}
\newcommand{\PeCfIkSZero}{19.25}
\newcommand{\PeCfIkVsLcpp}{19.6}
\newcommand{\PeCfIkXxsZero}{20.15}
\newcommand{\PeCfLcppDeep}{21.11}
\newcommand{\PeCfLcppMtpDeep}{5.98}
\newcommand{\PeCfLcppMtpZero}{6.43}
\newcommand{\PeCfLcppZero}{25.06}
\newcommand{\PeCfNifSBelowXxs}{6.5}
\newcommand{\PeCfNifSDeep}{38.24}
\newcommand{\PeCfNifSZero}{42.93}
\newcommand{\PeCfNifXxsDeep}{40.64}
\newcommand{\PeCfNifXxsZero}{45.91}
\newcommand{\PeCfPNifSDeep}{15.7}
\newcommand{\PeCfPNifSDeepHi}{15.9}
\newcommand{\PeCfPNifSDeepLo}{15.4}
\newcommand{\PeCfPNifSZero}{11.5}
\newcommand{\PeCfPNifSZeroHi}{12.2}
\newcommand{\PeCfPNifSZeroLo}{11.1}
\newcommand{\PeCfPNifXxsDeep}{8.9}
\newcommand{\PeCfPNifXxsDeepHi}{9.1}
\newcommand{\PeCfPNifXxsDeepLo}{8.6}
\newcommand{\PeCfPNifXxsZero}{4.3}
\newcommand{\PeCfPNifXxsZeroHi}{4.9}
\newcommand{\PeCfPNifXxsZeroLo}{3.9}
\newcommand{\PeCfPRafDeep}{6.5}
\newcommand{\PeCfPRafDeepEq}{6.1}
\newcommand{\PeCfPRafDeepHi}{11.8}
\newcommand{\PeCfPRafDeepLo}{2.1}
\newcommand{\PeCfPRafZero}{4.6}
\newcommand{\PeCfPRafZeroEq}{4.4}
\newcommand{\PeCfPRafZeroHi}{6.7}
\newcommand{\PeCfPRafZeroLo}{1.2}
\newcommand{\PeCfRafDeep}{63.04}
\newcommand{\PeCfRafVsFtPct}{3.7}
\newcommand{\PeCfRafZero}{68.61}
\newcommand{\PeCfSixMax}{1.6}
\newcommand{\PeCfSixMtpZero}{48.35}
\newcommand{\PeCfSixNdZero}{25.28}
\newcommand{\PeCfStockMtpDeep}{36.07}
\newcommand{\PeCfStockMtpZero}{47.59}
\newcommand{\PeCfStockNdDeep}{20.40}
\newcommand{\PeCfStockNdZero}{25.39}
\newcommand{\PeCfStopMax}{2.3}
\newcommand{\PeCfTimingCloseMs}{0.1}
\newcommand{\PeCfTimingCloseN}{57}
\newcommand{\PeCfTimingMaxMs}{47}
\newcommand{\PeCfTimingN}{114}
\newcommand{\PeCfWMtpDeep}{67.12}
\newcommand{\PeCfWMtpZero}{71.80}
\newcommand{\PeCfWNdDeep}{44.25}
\newcommand{\PeCfWNdZero}{47.87}
\newcommand{\PeCfWSixMtpZero}{72.16}
\newcommand{\PeCfWSixNdZero}{48.04}
\newcommand{\PeCfXBuunMtpZero}{1.81}
\newcommand{\PeCfXBuunNdZero}{1.82}
\newcommand{\PeCfXIkSZero}{2.49}
\newcommand{\PeCfXIkXxsZero}{2.38}
\newcommand{\PeCfXLcppDeep}{2.10}
\newcommand{\PeCfXLcppMtpDeep}{10.9}
\newcommand{\PeCfXLcppMtpZero}{11.5}
\newcommand{\PeCfXLcppZero}{1.91}
\newcommand{\PeCfXStockMtpDeep}{1.86}
\newcommand{\PeCfXStockMtpZero}{1.51}
\newcommand{\PeCfXStockNdDeep}{2.17}
\newcommand{\PeCfXStockNdZero}{1.89}
\newcommand{\PeCgAttnCut}{69}
\newcommand{\PeCgGdnCut}{17}
\newcommand{\PeCgMatrixCut}{48}
\newcommand{\PeCgMsCut}{49}
\newcommand{\PeCgSmallCut}{5}
\newcommand{\PeCgStockMtpAttnShare}{5.4}
\newcommand{\PeCgStockMtpAttnUs}{1,264}
\newcommand{\PeCgStockMtpGdnShare}{2.3}
\newcommand{\PeCgStockMtpGdnUs}{549}
\newcommand{\PeCgStockMtpMatrixShare}{88.3}
\newcommand{\PeCgStockMtpMatrixUs}{20,753}
\newcommand{\PeCgStockMtpMs}{26.28}
\newcommand{\PeCgStockMtpSmallShare}{4.0}
\newcommand{\PeCgStockMtpSmallUs}{949}
\newcommand{\PeCgStockNdAttnShare}{3.1}
\newcommand{\PeCgStockNdAttnUs}{1,344}
\newcommand{\PeCgStockNdGdnShare}{1.6}
\newcommand{\PeCgStockNdGdnUs}{692}
\newcommand{\PeCgStockNdMatrixShare}{91.9}
\newcommand{\PeCgStockNdMatrixUs}{40,322}
\newcommand{\PeCgStockNdMs}{45.58}
\newcommand{\PeCgStockNdSmallShare}{3.5}
\newcommand{\PeCgStockNdSmallUs}{1,514}
\newcommand{\PeCgWAttnShare}{3.1}
\newcommand{\PeCgWAttnUs}{390}
\newcommand{\PeCgWDeepAttnShare}{11.0}
\newcommand{\PeCgWDeepAttnUs}{1,508}
\newcommand{\PeCgWDeepGdnShare}{3.4}
\newcommand{\PeCgWDeepGdnUs}{462}
\newcommand{\PeCgWDeepMatrixShare}{78.9}
\newcommand{\PeCgWDeepMatrixUs}{10,823}
\newcommand{\PeCgWDeepMs}{14.51}
\newcommand{\PeCgWDeepSmallShare}{6.7}
\newcommand{\PeCgWDeepSmallUs}{918}
\newcommand{\PeCgWGdnShare}{3.6}
\newcommand{\PeCgWGdnUs}{454}
\newcommand{\PeCgWMatrixShare}{86.0}
\newcommand{\PeCgWMatrixUs}{10,695}
\newcommand{\PeCgWMs}{13.52}
\newcommand{\PeCgWSmallShare}{7.3}
\newcommand{\PeCgWSmallUs}{904}
\newcommand{\PeChunk}{1,024}
\newcommand{\PeCliffArms}{5}
\newcommand{\PeCliffArmsWord}{five}
\newcommand{\PeCliffVtMax}{1}
\newcommand{\PeCrossDiff}{\ensuremath{+5.2}}
\newcommand{\PeCrossHi}{\ensuremath{+14.6}}
\newcommand{\PeCrossLo}{\ensuremath{-4.2}}
\newcommand{\PeCrossOtt}{12}
\newcommand{\PeCrossP}{0.46}
\newcommand{\PeCrossTen}{17}
\newcommand{\PeCuda}{12.8}
\newcommand{\PeCycleEEMs}{38.7}
\newcommand{\PeCycleFFMs}{57.5}
\newcommand{\PeCycleFpMs}{297}
\newcommand{\PeDecodeRatio}{1.8}
\newcommand{\PeDepthAfter}{25.35}
\newcommand{\PeDepthBefore}{32.75}
\newcommand{\PeDepthCallRatio}{5.2}
\newcommand{\PeDepthChange}{22.6}
\newcommand{\PeDepthFdCall}{135}
\newcommand{\PeDepthFdUs}{2,149}
\newcommand{\PeDepthTritonCall}{699}
\newcommand{\PeDepthTritonUs}{11,140}
\newcommand{\PeDepthTwoFOne}{28.89}
\newcommand{\PeDepthTwoGain}{37.5}
\newcommand{\PeDepthTwoOff}{2.6}
\newcommand{\PeDepthTwoW}{39.73}
\newcommand{\PeDepthTwoZeroGain}{4.2}
\newcommand{\PeDfaMtpGain}{1.52}
\newcommand{\PeDfaNifXxsZero}{46.01}
\newcommand{\PeDfaPNifXxsZero}{4.2}
\newcommand{\PeDfaWMtpDeep}{66.48}
\newcommand{\PeDfaWMtpZero}{73.00}
\newcommand{\PeDfaWNdZero}{47.96}
\newcommand{\PeDfbAccDeepCode}{51}
\newcommand{\PeDfbAccDeepMath}{65}
\newcommand{\PeDfbAccDeepProse}{17}
\newcommand{\PeDfbAccZeroCode}{69}
\newcommand{\PeDfbAccZeroMath}{74}
\newcommand{\PeDfbAccZeroProse}{17}
\newcommand{\PeDfbFastestDeep}{holds}
\newcommand{\PeDfbFastestZero}{holds}
\newcommand{\PeDfbFtwoMax}{1.4}
\newcommand{\PeDfbMtpGain}{1.52}
\newcommand{\PeDfbNifCtx}{22,144}
\newcommand{\PeDfbNifDeepCtx}{29,696}
\newcommand{\PeDfbNifDeepRung}{6 of 6}
\newcommand{\PeDfbNifDfDeep}{48.90}
\newcommand{\PeDfbNifDfGain}{1.25}
\newcommand{\PeDfbNifDfZero}{57.52}
\newcommand{\PeDfbNifXxsZero}{45.92}
\newcommand{\PeDfbPNifDfDeep}{34.2}
\newcommand{\PeDfbPNifDfDeepHi}{44.3}
\newcommand{\PeDfbPNifDfDeepLo}{23.9}
\newcommand{\PeDfbPNifDfDeepWins}{3}
\newcommand{\PeDfbPNifDfZero}{26.1}
\newcommand{\PeDfbPNifDfZeroHi}{28.6}
\newcommand{\PeDfbPNifDfZeroLo}{23.4}
\newcommand{\PeDfbPNifDfZeroWins}{3}
\newcommand{\PeDfbPpNifDfDeepCode}{64.5}
\newcommand{\PeDfbPpNifDfDeepMath}{76.5}
\newcommand{\PeDfbPpNifDfDeepProse}{30.5}
\newcommand{\PeDfbPpNifDfZeroCode}{88.1}
\newcommand{\PeDfbPpNifDfZeroMath}{93.7}
\newcommand{\PeDfbPpNifDfZeroProse}{33.2}
\newcommand{\PeDfbPpNifXxsZeroCode}{45.9}
\newcommand{\PeDfbPpNifXxsZeroMath}{45.9}
\newcommand{\PeDfbPpNifXxsZeroProse}{45.9}
\newcommand{\PeDfbPpWMtpDeepCode}{79.2}
\newcommand{\PeDfbPpWMtpDeepMath}{85.4}
\newcommand{\PeDfbPpWMtpDeepProse}{47.9}
\newcommand{\PeDfbPpWMtpZeroCode}{87.8}
\newcommand{\PeDfbPpWMtpZeroMath}{91.8}
\newcommand{\PeDfbPpWMtpZeroProse}{52.4}
\newcommand{\PeDfbWMtpDeep}{65.61}
\newcommand{\PeDfbWMtpZero}{72.50}
\newcommand{\PeDfbWNdZero}{47.80}
\newcommand{\PeDraftTokens}{6}
\newcommand{\PeDriver}{580.178.04}
\newcommand{\PeEEEightyOver}{25}
\newcommand{\PeEESixtyFourFree}{1,167}
\newcommand{\PeEmbedGiB}{2.37}
\newcommand{\PeExlVersionShort}{1.6.5}
\newcommand{\PeExtraEE}{91}
\newcommand{\PeExtraEF}{117}
\newcommand{\PeExtraFF}{108}
\newcommand{\PeFileGbIxxs}{7.27}
\newcommand{\PeFilesExlBpw}{2.94}
\newcommand{\PeFilesExlParams}{27.78}
\newcommand{\PeFilesExlTrellisBpw}{2.05}
\newcommand{\PeFilesExlTrellisParams}{26.02}
\newcommand{\PeFilesGgufParams}{26.90}
\newcommand{\PeFillGapMax}{562}
\newcommand{\PeFillGapMin}{524}
\newcommand{\PeFleetAssessedPredictions}{11}
\newcommand{\PeFleetCardAGsmMtp}{488}
\newcommand{\PeFleetCardAGsmNd}{486}
\newcommand{\PeFleetCardAMmluMtp}{434}
\newcommand{\PeFleetCardAMmluNd}{431}
\newcommand{\PeFleetCardAMtpDeep}{98.391}
\newcommand{\PeFleetCardAMtpZero}{109.166}
\newcommand{\PeFleetCardANdDeep}{58.714}
\newcommand{\PeFleetCardANdZero}{64.193}
\newcommand{\PeFleetCardBGsmMtp}{488}
\newcommand{\PeFleetCardBGsmNd}{485}
\newcommand{\PeFleetCardBMmluMtp}{431}
\newcommand{\PeFleetCardBMmluNd}{436}
\newcommand{\PeFleetCardBMtpDeep}{60.548}
\newcommand{\PeFleetCardBMtpZero}{69.174}
\newcommand{\PeFleetCardBNdDeep}{40.951}
\newcommand{\PeFleetCardBNdZero}{45.774}
\newcommand{\PeFleetQualityN}{500}
\newcommand{\PeFleetSpeedRequestsPerCell}{9}
\newcommand{\PeFoneAdoptTwentyEight}{40.01}
\newcommand{\PeFoneRatioMtpTwentyEight}{1.38}
\newcommand{\PeFoneRatioTwentyEight}{2.13}
\newcommand{\PeFoneStockTwentyEight}{18.78}
\newcommand{\PeFreshComplete}{9}
\newcommand{\PeFreshEligible}{4}
\newcommand{\PeFreshLcppMtpDeep}{5.80}
\newcommand{\PeFreshLcppMtpZero}{6.31}
\newcommand{\PeFreshLcppNdDeep}{21.11}
\newcommand{\PeFreshLcppNdZero}{24.96}
\newcommand{\PeFreshNifDfDeep}{51.62}
\newcommand{\PeFreshNifDfDeepCode}{64.51}
\newcommand{\PeFreshNifDfDeepMath}{87.84}
\newcommand{\PeFreshNifDfDeepProse}{32.02}
\newcommand{\PeFreshNifDfZero}{54.37}
\newcommand{\PeFreshNifDfZeroCode}{60.34}
\newcommand{\PeFreshNifDfZeroMath}{94.85}
\newcommand{\PeFreshNifDfZeroProse}{35.64}
\newcommand{\PeFreshNifNdDeep}{41.38}
\newcommand{\PeFreshNifNdZero}{47.23}
\newcommand{\PeFreshPLcppMtpDeep}{1050.0}
\newcommand{\PeFreshPLcppMtpDeepHi}{1082.0}
\newcommand{\PeFreshPLcppMtpDeepLo}{1011.0}
\newcommand{\PeFreshPLcppMtpZero}{1091.1}
\newcommand{\PeFreshPLcppMtpZeroHi}{1121.3}
\newcommand{\PeFreshPLcppMtpZeroLo}{1067.5}
\newcommand{\PeFreshPLcppNdDeep}{109.06}
\newcommand{\PeFreshPLcppNdDeepHi}{109.30}
\newcommand{\PeFreshPLcppNdDeepLo}{108.64}
\newcommand{\PeFreshPLcppNdZero}{91.25}
\newcommand{\PeFreshPLcppNdZeroHi}{92.11}
\newcommand{\PeFreshPLcppNdZeroLo}{90.51}
\newcommand{\PeFreshPNifDfDeep}{28.4}
\newcommand{\PeFreshPNifDfDeepHi}{37.7}
\newcommand{\PeFreshPNifDfDeepLo}{21.6}
\newcommand{\PeFreshPNifDfZero}{33.7}
\newcommand{\PeFreshPNifDfZeroHi}{45.6}
\newcommand{\PeFreshPNifDfZeroLo}{24.8}
\newcommand{\PeFreshPNifNdDeep}{6.80}
\newcommand{\PeFreshPNifNdDeepHi}{7.17}
\newcommand{\PeFreshPNifNdDeepLo}{6.46}
\newcommand{\PeFreshPNifNdZero}{1.09}
\newcommand{\PeFreshPNifNdZeroHi}{1.24}
\newcommand{\PeFreshPNifNdZeroLo}{0.95}
\newcommand{\PeFreshPStockMtpZero}{48.9}
\newcommand{\PeFreshPStockMtpZeroHi}{50.4}
\newcommand{\PeFreshPStockMtpZeroLo}{47.7}
\newcommand{\PeFreshPlanned}{14}
\newcommand{\PeFreshStockMtpZero}{48.52}
\newcommand{\PeFreshWLcppMtpDeep}{66.73}
\newcommand{\PeFreshWLcppMtpZero}{75.11}
\newcommand{\PeFreshWLcppNdDeep}{44.13}
\newcommand{\PeFreshWLcppNdZero}{47.73}
\newcommand{\PeFreshWNifDfDeep}{66.26}
\newcommand{\PeFreshWNifDfDeepCode}{82.62}
\newcommand{\PeFreshWNifDfDeepMath}{86.22}
\newcommand{\PeFreshWNifDfDeepProse}{47.33}
\newcommand{\PeFreshWNifDfZero}{72.70}
\newcommand{\PeFreshWNifDfZeroCode}{88.54}
\newcommand{\PeFreshWNifDfZeroMath}{92.12}
\newcommand{\PeFreshWNifDfZeroProse}{52.32}
\newcommand{\PeFreshWNifNdDeep}{44.19}
\newcommand{\PeFreshWNifNdZero}{47.75}
\newcommand{\PeFreshWStockMtpZero}{72.23}
\newcommand{\PeFreshXLcppMtpDeep}{11.50}
\newcommand{\PeFreshXLcppMtpZero}{11.91}
\newcommand{\PeFreshXLcppNdDeep}{2.09}
\newcommand{\PeFreshXLcppNdZero}{1.91}
\newcommand{\PeFreshXNifDfDeep}{1.28}
\newcommand{\PeFreshXNifDfZero}{1.34}
\newcommand{\PeFreshXNifNdDeep}{1.07}
\newcommand{\PeFreshXNifNdZero}{1.01}
\newcommand{\PeFreshXStockMtpZero}{1.49}
\newcommand{\PeFtAcceptHi}{0.56}
\newcommand{\PeFtAcceptLo}{0.53}
\newcommand{\PeFtAdoptMtpTwentyEight}{50.51}
\newcommand{\PeFtAdoptMtpZero}{54.60}
\newcommand{\PeFtAdoptTwentyEight}{41.91}
\newcommand{\PeFtAdoptZero}{45.45}
\newcommand{\PeFtFirst}{14:09}
\newcommand{\PeFtLast}{14:51}
\newcommand{\PeFtLoadEnd}{10.9}
\newcommand{\PeFtLoadStart}{3.2}
\newcommand{\PeFtModeTwoOverAdoptPct}{31.2}
\newcommand{\PeFtMtpCutPct}{35}
\newcommand{\PeFtMtpOverStockHi}{+61.7}
\newcommand{\PeFtMtpOverStockLo}{+47.4}
\newcommand{\PeFtMtpOverStockPct}{54.3}
\newcommand{\PeFtMtpOverStockWins}{3}
\newcommand{\PeFtMtpRatio}{1.54}
\newcommand{\PeFtOverRafTwentyEightEqPct}{5.7}
\newcommand{\PeFtOverRafTwentyEightHi}{+9.2}
\newcommand{\PeFtOverRafTwentyEightLo}{+1.0}
\newcommand{\PeFtOverRafTwentyEightPct}{5.5}
\newcommand{\PeFtOverRafTwentyEightWins}{3}
\newcommand{\PeFtOverRafZeroEqPct}{7.1}
\newcommand{\PeFtOverRafZeroHi}{+11.7}
\newcommand{\PeFtOverRafZeroLo}{+5.2}
\newcommand{\PeFtOverRafZeroPct}{8.2}
\newcommand{\PeFtOverRafZeroWins}{3}
\newcommand{\PeFtPpSTwoCode}{87.5}
\newcommand{\PeFtPpSTwoMath}{92.5}
\newcommand{\PeFtPpSTwoProse}{50.8}
\newcommand{\PeFtRafMtpTwentyEight}{63.09}
\newcommand{\PeFtRafMtpZero}{66.18}
\newcommand{\PeFtRows}{126}
\newcommand{\PeFtSOneMtpTwentyEight}{61.33}
\newcommand{\PeFtSOneMtpZero}{65.50}
\newcommand{\PeFtSOneOverRafZeroPct}{1.0}
\newcommand{\PeFtSOneOverRafZeroWins}{0}
\newcommand{\PeFtSTwoMtpTwentyEight}{66.57}
\newcommand{\PeFtSTwoMtpZero}{71.61}
\newcommand{\PeFtSplitKMtpPct}{20.0}
\newcommand{\PeFtSplitKMtpTwentyEightPct}{21.4}
\newcommand{\PeFtSplitKMtpTwentyEightWins}{3}
\newcommand{\PeFtSplitKMtpWins}{3}
\newcommand{\PeFtSplitKUndraftedPct}{4.6}
\newcommand{\PeFtSplitKUndraftedTwentyEightPct}{4.7}
\newcommand{\PeFtSplitKUndraftedTwentyEightWins}{3}
\newcommand{\PeFtSplitKUndraftedWins}{3}
\newcommand{\PeFtStockMtpZero}{46.40}
\newcommand{\PeFtStockZero}{25.33}
\newcommand{\PeFtTempHi}{78}
\newcommand{\PeFtTempLo}{69}
\newcommand{\PeFtUndraftedOverStockHi}{+88.6}
\newcommand{\PeFtUndraftedOverStockLo}{+86.1}
\newcommand{\PeFtUndraftedOverStockPct}{87.6}
\newcommand{\PeFtUndraftedOverStockWins}{3}
\newcommand{\PeFtUndraftedRatio}{1.88}
\newcommand{\PeFtWMs}{21.0}
\newcommand{\PeFtWTwentyEight}{43.88}
\newcommand{\PeFtWZero}{47.53}
\newcommand{\PeGapAbs}{27.1}
\newcommand{\PeGapAbsRound}{27}
\newcommand{\PeGapDiff}{\ensuremath{-27.1}}
\newcommand{\PeGapEEAbs}{14.6}
\newcommand{\PeGapEEAbsRound}{15}
\newcommand{\PeGapEEDiff}{\ensuremath{-14.6}}
\newcommand{\PeGapEEHi}{\ensuremath{0.0}}
\newcommand{\PeGapEEK}{35}
\newcommand{\PeGapEELo}{\ensuremath{-29.2}}
\newcommand{\PeGapEEOtt}{14}
\newcommand{\PeGapEEP}{0.19}
\newcommand{\PeGapEETen}{7}
\newcommand{\PeGapFpAbs}{18.8}
\newcommand{\PeGapFpAbsRound}{19}
\newcommand{\PeGapFpDiff}{\ensuremath{-18.8}}
\newcommand{\PeGapFpHi}{\ensuremath{-4.2}}
\newcommand{\PeGapFpK}{37}
\newcommand{\PeGapFpLo}{\ensuremath{-33.3}}
\newcommand{\PeGapFpOtt}{14}
\newcommand{\PeGapFpP}{0.064}
\newcommand{\PeGapFpTen}{5}
\newcommand{\PeGapHi}{\ensuremath{-14.6}}
\newcommand{\PeGapLo}{\ensuremath{-41.7}}
\newcommand{\PeGapN}{48}
\newcommand{\PeGapNinetySixK}{28}
\newcommand{\PeGapOtt}{16}
\newcommand{\PeGapP}{0.0044}
\newcommand{\PeGapSixtyFourAcc}{85.4}
\newcommand{\PeGapSixtyFourK}{41}
\newcommand{\PeGapTen}{3}
\newcommand{\PeGbsNif}{299}
\newcommand{\PeGbsW}{315}
\newcommand{\PeHOneAbs}{0.0}
\newcommand{\PeHOneDiff}{\ensuremath{0.0}}
\newcommand{\PeHOneHi}{\ensuremath{0.0}}
\newcommand{\PeHOneLo}{\ensuremath{0.0}}
\newcommand{\PeHOneN}{96}
\newcommand{\PeHOneOtt}{0}
\newcommand{\PeHOneP}{1}
\newcommand{\PeHOneSameTexts}{94}
\newcommand{\PeHOneTen}{0}
\newcommand{\PeHThreeAbs}{3.1}
\newcommand{\PeHThreeDiff}{\ensuremath{-3.1}}
\newcommand{\PeHThreeHi}{\ensuremath{0.0}}
\newcommand{\PeHThreeHolm}{0.75}
\newcommand{\PeHThreeLo}{\ensuremath{-7.3}}
\newcommand{\PeHThreeN}{96}
\newcommand{\PeHThreeOtt}{4}
\newcommand{\PeHThreeP}{0.375}
\newcommand{\PeHThreeTen}{1}
\newcommand{\PeHTwoAbs}{3.1}
\newcommand{\PeHTwoDiff}{\ensuremath{-3.1}}
\newcommand{\PeHTwoDisc}{23}
\newcommand{\PeHTwoDiscPct}{24.0}
\newcommand{\PeHTwoHi}{\ensuremath{+5.2}}
\newcommand{\PeHTwoLo}{\ensuremath{-11.5}}
\newcommand{\PeHTwoN}{96}
\newcommand{\PeHTwoOtt}{13}
\newcommand{\PeHTwoP}{0.678}
\newcommand{\PeHTwoSAbs}{11.5}
\newcommand{\PeHTwoSDiff}{\ensuremath{+11.5}}
\newcommand{\PeHTwoSDisc}{17}
\newcommand{\PeHTwoSDiscPct}{17.7}
\newcommand{\PeHTwoSHi}{\ensuremath{+18.8}}
\newcommand{\PeHTwoSLo}{\ensuremath{+5.2}}
\newcommand{\PeHTwoSN}{96}
\newcommand{\PeHTwoSOtt}{3}
\newcommand{\PeHTwoSP}{0.0127}
\newcommand{\PeHTwoSSameK}{6}
\newcommand{\PeHTwoSSamePct}{37.5}
\newcommand{\PeHTwoSTen}{14}
\newcommand{\PeHTwoTen}{10}
\newcommand{\PeHeadroomMiB}{512}
\newcommand{\PeIntEightAfter}{27.78}
\newcommand{\PeIntEightAlone}{22.9}
\newcommand{\PeIntEightBefore}{45.58}
\newcommand{\PeIntEightChange}{39}
\newcommand{\PeIntEightUnfused}{6.8}
\newcommand{\PeItemSDiffPct}{0.9}
\newcommand{\PeKOneAbs}{0.0}
\newcommand{\PeKOneBound}{3.1}
\newcommand{\PeKOneChanged}{2}
\newcommand{\PeKOneDiff}{\ensuremath{0.0}}
\newcommand{\PeKOneDiffer}{3}
\newcommand{\PeKOneHi}{\ensuremath{+2.1}}
\newcommand{\PeKOneLo}{\ensuremath{-2.1}}
\newcommand{\PeKOneN}{96}
\newcommand{\PeKOneOtt}{1}
\newcommand{\PeKOneP}{1}
\newcommand{\PeKOneSame}{93}
\newcommand{\PeKOneTen}{1}
\newcommand{\PeKTwoAbs}{11.5}
\newcommand{\PeKTwoChanged}{17}
\newcommand{\PeKTwoDiff}{\ensuremath{+11.5}}
\newcommand{\PeKTwoDiffer}{30}
\newcommand{\PeKTwoHi}{\ensuremath{+18.8}}
\newcommand{\PeKTwoLo}{\ensuremath{+5.2}}
\newcommand{\PeKTwoN}{96}
\newcommand{\PeKTwoOtt}{3}
\newcommand{\PeKTwoP}{0.0127}
\newcommand{\PeKTwoSame}{66}
\newcommand{\PeKTwoTen}{14}
\newcommand{\PeKlFortyEight}{\ensuremath{3.50\times10^{-4}}}
\newcommand{\PeKlMax}{0.024}
\newcommand{\PeKlMinutes}{39.1}
\newcommand{\PeKlPerLength}{288}
\newcommand{\PeKlPositions}{864}
\newcommand{\PeKlSixtyFour}{\ensuremath{2.56\times10^{-5}}}
\newcommand{\PeKlTwentyEight}{\ensuremath{1.14\times10^{-5}}}
\newcommand{\PeLimitMiB}{10,752}
\newcommand{\PeMaxTemp}{80}
\newcommand{\PeNConfigs}{16}
\newcommand{\PeNRequests}{1,048}
\newcommand{\PeNSessions}{16}
\newcommand{\PeNStarts}{16}
\newcommand{\PeNeedles}{12}
\newcommand{\PeNifDeepDflash}{56.5}
\newcommand{\PeNifDeepTokens}{22,224}
\newcommand{\PeNifDeepUndrafted}{40.4}
\newcommand{\PeNifDflash}{52.1}
\newcommand{\PeNifLeadPct}{9}
\newcommand{\PeNifUndrafted}{43.7}
\newcommand{\PeNineFixedTwo}{8.1}
\newcommand{\PeNineFixedTwoHalf}{0.2}
\newcommand{\PeNineFourBitGbs}{567}
\newcommand{\PeNineFuseGateUp}{0.36}
\newcommand{\PeNineFuseQkv}{0.94}
\newcommand{\PeNineFuseQkvZ}{0.92}
\newcommand{\PeNineIpwOne}{1.26}
\newcommand{\PeNineIpwTwo}{1.21}
\newcommand{\PeNineIssuePct}{17}
\newcommand{\PeNineOneTwoHi}{97}
\newcommand{\PeNineOneTwoKv}{125}
\newcommand{\PeNineOneTwoLo}{95}
\newcommand{\PeNineOneTwoPerWeight}{93}
\newcommand{\PeNinePerWeightOneX}{2.2}
\newcommand{\PeNinePerWeightTwoPct}{16}
\newcommand{\PeNinePhEpi}{2.4}
\newcommand{\PeNinePhFull}{7.5}
\newcommand{\PeNinePhLoop}{1.7}
\newcommand{\PeNinePhNull}{0.8}
\newcommand{\PeNinePhStage}{1.0}
\newcommand{\PeNineShare}{61}
\newcommand{\PeNineShareRated}{50}
\newcommand{\PeNinetySixUnderLimit}{103}
\newcommand{\PeNoiseHi}{19.4}
\newcommand{\PeNoiseK}{0}
\newcommand{\PeNoiseN}{16}
\newcommand{\PeNoiseSameK}{4}
\newcommand{\PeNzAbs}{1.0}
\newcommand{\PeNzAccA}{71.9}
\newcommand{\PeNzAccB}{70.8}
\newcommand{\PeNzDiff}{\ensuremath{+1.0}}
\newcommand{\PeNzDisc}{3}
\newcommand{\PeNzFirstSixteen}{2}
\newcommand{\PeNzHi}{\ensuremath{+4.2}}
\newcommand{\PeNzKTwoDisc}{17.7}
\newcommand{\PeNzLo}{\ensuremath{-2.1}}
\newcommand{\PeNzN}{96}
\newcommand{\PeNzOtt}{1}
\newcommand{\PeNzP}{1}
\newcommand{\PeNzSame}{92}
\newcommand{\PeNzTen}{2}
\newcommand{\PeNzWilsonHi}{8.8}
\newcommand{\PeOnCardGiB}{6.27}
\newcommand{\PePOneMax}{222}
\newcommand{\PePOneMin}{204}
\newcommand{\PePeakEEEighty}{10,777}
\newcommand{\PePeakEESeventyTwo}{10,477}
\newcommand{\PePeakEESixtyFour}{10,097}
\newcommand{\PePeakEFEightyEight}{10,335}
\newcommand{\PePeakEFNinetySix}{10,555}
\newcommand{\PePeakEFSixtyFour}{9,555}
\newcommand{\PePeakFFNinetySix}{9,735}
\newcommand{\PePeakFFOneThirtySix}{10,495}
\newcommand{\PePeakFFOneTwentyEight}{10,355}
\newcommand{\PePeakFFSixtyFour}{9,015}
\newcommand{\PePeakFpForty}{10,645}
\newcommand{\PePeakFpThirtySix}{10,365}
\newcommand{\PePeakFpThirtyTwo}{10,105}
\newcommand{\PePeakRss}{5.9}
\newcommand{\PePrefillArmMax}{575}
\newcommand{\PePrefillArmMin}{571}
\newcommand{\PePrefillEEEighty}{148}
\newcommand{\PePrefillEESeventyTwo}{131}
\newcommand{\PePrefillEESixtyFour}{108}
\newcommand{\PePrefillEFEightyEight}{167}
\newcommand{\PePrefillEFNinetySix}{187}
\newcommand{\PePrefillEFSixtyFour}{113}
\newcommand{\PePrefillFFNinetySix}{187}
\newcommand{\PePrefillFFOneThirtySix}{294}
\newcommand{\PePrefillFFOneTwentyEight}{271}
\newcommand{\PePrefillFFSixtyFour}{113}
\newcommand{\PePrefillFpForty}{65}
\newcommand{\PePrefillFpThirtySix}{58}
\newcommand{\PePrefillFpThirtyTwo}{51}
\newcommand{\PePricedMax}{9.2}
\newcommand{\PePricedMin}{3.1}
\newcommand{\PeProfFinPct}{2}
\newcommand{\PeProfStockMtpPct}{22}
\newcommand{\PeProfStockPct}{15}
\newcommand{\PeProfWMtpPct}{3}
\newcommand{\PeProjUnderMax}{222}
\newcommand{\PeProjUnderMin}{62}
\newcommand{\PePromptEEEighty}{81,395}
\newcommand{\PePromptEESeventyTwo}{73,179}
\newcommand{\PePromptEESixtyFour}{64,980}
\newcommand{\PePromptEFEightyEight}{89,588}
\newcommand{\PePromptEFNinetySix}{97,774}
\newcommand{\PePromptEFSixtyFour}{64,980}
\newcommand{\PePromptFFNinetySix}{97,774}
\newcommand{\PePromptFFOneThirtySix}{138,702}
\newcommand{\PePromptFFOneTwentyEight}{130,546}
\newcommand{\PePromptFFSixtyFour}{64,980}
\newcommand{\PePromptFpForty}{40,432}
\newcommand{\PePromptFpThirtySix}{36,337}
\newcommand{\PePromptFpThirtyTwo}{32,215}
\newcommand{\PePromptMaxNinetySix}{97,711}
\newcommand{\PePromptMaxSixtyFour}{64,943}
\newcommand{\PePromptMaxThirtyTwo}{32,175}
\newcommand{\PePromptMinNinetySix}{97,663}
\newcommand{\PePromptMinSixtyFour}{64,907}
\newcommand{\PePromptMinThirtyTwo}{32,128}
\newcommand{\PePython}{3.12.16}
\newcommand{\PeQGsmExlMtpAcc}{97.6}
\newcommand{\PeQGsmExlMtpTok}{327}
\newcommand{\PeQGsmExlMtpTps}{53.4}
\newcommand{\PeQGsmExlMtpWall}{6.45}
\newcommand{\PeQGsmExlNoAcc}{97.6}
\newcommand{\PeQGsmExlNoTok}{326}
\newcommand{\PeQGsmExlNoTps}{23.6}
\newcommand{\PeQGsmExlNoWall}{14.22}
\newcommand{\PeQGsmExlWAcc}{97.6}
\newcommand{\PeQGsmExlWTok}{326}
\newcommand{\PeQGsmExlWTps}{85.9}
\newcommand{\PeQGsmExlWWall}{4.07}
\newcommand{\PeQGsmIsAcc}{96.6}
\newcommand{\PeQGsmIsMtpAcc}{97.4}
\newcommand{\PeQGsmIsMtpTok}{368}
\newcommand{\PeQGsmIsMtpTps}{48.7}
\newcommand{\PeQGsmIsMtpWall}{8.17}
\newcommand{\PeQGsmIsTok}{364}
\newcommand{\PeQGsmIsTps}{23.4}
\newcommand{\PeQGsmIsWall}{16.01}
\newcommand{\PeQGsmIxxsAcc}{96.4}
\newcommand{\PeQGsmIxxsDfAcc}{96.2}
\newcommand{\PeQGsmIxxsDfTok}{376}
\newcommand{\PeQGsmIxxsDfTps}{54.2}
\newcommand{\PeQGsmIxxsDfWall}{7.69}
\newcommand{\PeQGsmIxxsTok}{376}
\newcommand{\PeQGsmIxxsTps}{24.7}
\newcommand{\PeQGsmIxxsWall}{15.69}
\newcommand{\PeQGsmMax}{97.6}
\newcommand{\PeQGsmMin}{96.2}
\newcommand{\PeQLongerMax}{16}
\newcommand{\PeQLongerMin}{12}
\newcommand{\PeQMmluExlMtpAcc}{86.8}
\newcommand{\PeQMmluExlNoAcc}{86.6}
\newcommand{\PeQMmluIsAcc}{86.2}
\newcommand{\PeQMmluIsMtpAcc}{86.0}
\newcommand{\PeQMmluIxxsAcc}{82.8}
\newcommand{\PeQMmluIxxsDfAcc}{83.2}
\newcommand{\PeQMtpGsmOtt}{0}
\newcommand{\PeQMtpGsmTen}{0}
\newcommand{\PeQMtpMmluOtt}{3}
\newcommand{\PeQMtpMmluTen}{2}
\newcommand{\PeQPGsmIsAbs}{1.0}
\newcommand{\PeQPGsmIsDiff}{\ensuremath{-1.0}}
\newcommand{\PeQPGsmIsHolm}{1.00}
\newcommand{\PeQPGsmIsOtt}{3}
\newcommand{\PeQPGsmIsP}{0.23}
\newcommand{\PeQPGsmIsTen}{8}
\newcommand{\PeQPGsmIxxsAbs}{1.2}
\newcommand{\PeQPGsmIxxsDiff}{\ensuremath{-1.2}}
\newcommand{\PeQPGsmIxxsHolm}{1.00}
\newcommand{\PeQPGsmIxxsOtt}{3}
\newcommand{\PeQPGsmIxxsP}{0.15}
\newcommand{\PeQPGsmIxxsTen}{9}
\newcommand{\PeQPMmluDfAbs}{3.6}
\newcommand{\PeQPMmluDfDiff}{\ensuremath{-3.6}}
\newcommand{\PeQPMmluDfHolm}{0.23}
\newcommand{\PeQPMmluDfOtt}{17}
\newcommand{\PeQPMmluDfP}{0.018}
\newcommand{\PeQPMmluDfTen}{35}
\newcommand{\PeQPMmluIsAbs}{0.4}
\newcommand{\PeQPMmluIsDiff}{\ensuremath{-0.4}}
\newcommand{\PeQPMmluIsHolm}{1.00}
\newcommand{\PeQPMmluIsOtt}{16}
\newcommand{\PeQPMmluIsP}{0.86}
\newcommand{\PeQPMmluIsTen}{18}
\newcommand{\PeQPMmluIxxsAbs}{3.8}
\newcommand{\PeQPMmluIxxsDiff}{\ensuremath{-3.8}}
\newcommand{\PeQPMmluIxxsHolm}{0.23}
\newcommand{\PeQPMmluIxxsOtt}{19}
\newcommand{\PeQPMmluIxxsP}{0.016}
\newcommand{\PeQPMmluIxxsTen}{38}
\newcommand{\PeQRatioIsMtp}{2.01}
\newcommand{\PeQRatioIxxsDf}{1.89}
\newcommand{\PeQTokRatioIsMtp}{1.13}
\newcommand{\PeQTokRatioIxxsDf}{1.16}
\newcommand{\PeQTpsRatioIsMtp}{1.76}
\newcommand{\PeQTpsRatioIxxsDf}{1.58}
\newcommand{\PeQWElevenAcc}{96.8}
\newcommand{\PeQWElevenDiff}{\ensuremath{-0.8}}
\newcommand{\PeQWElevenMtpAcc}{97.6}
\newcommand{\PeQWElevenMtpTok}{324}
\newcommand{\PeQWElevenMtpTps}{63.2}
\newcommand{\PeQWElevenMtpWall}{5.44}
\newcommand{\PeQWElevenOtt}{0}
\newcommand{\PeQWElevenP}{0.125}
\newcommand{\PeQWElevenTen}{4}
\newcommand{\PeQWElevenTok}{328}
\newcommand{\PeQWElevenTps}{42.6}
\newcommand{\PeQWElevenWall}{7.95}
\newcommand{\PeQWIdent}{437}
\newcommand{\PeQWModeTwoRefAcc}{97.6}
\newcommand{\PeQWModeTwoRefSame}{499}
\newcommand{\PeQWSame}{500}
\newcommand{\PeROneAbs}{2.1}
\newcommand{\PeROneDiff}{\ensuremath{-2.1}}
\newcommand{\PeROneDiffer}{32}
\newcommand{\PeROneEightAcc}{64.6}
\newcommand{\PeROneEightHi}{73.4}
\newcommand{\PeROneEightK}{62}
\newcommand{\PeROneEightLo}{54.6}
\newcommand{\PeROneEightN}{96}
\newcommand{\PeROneHi}{\ensuremath{+4.2}}
\newcommand{\PeROneLo}{\ensuremath{-8.3}}
\newcommand{\PeROneN}{96}
\newcommand{\PeROneOtt}{8}
\newcommand{\PeROneP}{0.791}
\newcommand{\PeROneRequests}{192}
\newcommand{\PeROneSame}{64}
\newcommand{\PeROneSplitAcc}{62.5}
\newcommand{\PeROneSplitHi}{71.5}
\newcommand{\PeROneSplitK}{60}
\newcommand{\PeROneSplitLo}{52.5}
\newcommand{\PeROneSplitN}{96}
\newcommand{\PeROneTen}{6}
\newcommand{\PeRThreeAAbs}{3.1}
\newcommand{\PeRThreeADiff}{\ensuremath{+3.1}}
\newcommand{\PeRThreeAHi}{\ensuremath{+7.3}}
\newcommand{\PeRThreeALo}{\ensuremath{-1.0}}
\newcommand{\PeRThreeAN}{96}
\newcommand{\PeRThreeAOtt}{2}
\newcommand{\PeRThreeAP}{0.453}
\newcommand{\PeRThreeATen}{5}
\newcommand{\PeRThreeBAbs}{8.3}
\newcommand{\PeRThreeBDiff}{\ensuremath{+8.3}}
\newcommand{\PeRThreeBHi}{\ensuremath{+14.6}}
\newcommand{\PeRThreeBLo}{\ensuremath{+2.1}}
\newcommand{\PeRThreeBN}{96}
\newcommand{\PeRThreeBOtt}{3}
\newcommand{\PeRThreeBP}{0.0574}
\newcommand{\PeRThreeBTen}{11}
\newcommand{\PeRThreeEightAcc}{70.8}
\newcommand{\PeRThreeFourAcc}{82.3}
\newcommand{\PeRThreeSixAcc}{74.0}
\newcommand{\PeRTwoAbs}{0.0}
\newcommand{\PeRTwoAccA}{77.1}
\newcommand{\PeRTwoAccB}{77.1}
\newcommand{\PeRTwoCheck}{16}
\newcommand{\PeRTwoDiff}{\ensuremath{0.0}}
\newcommand{\PeRTwoHi}{\ensuremath{+5.2}}
\newcommand{\PeRTwoLo}{\ensuremath{-5.2}}
\newcommand{\PeRTwoN}{96}
\newcommand{\PeRTwoOtt}{4}
\newcommand{\PeRTwoP}{1}
\newcommand{\PeRTwoTen}{4}
\newcommand{\PeReachEE}{72K}
\newcommand{\PeReachEF}{96K}
\newcommand{\PeReachFF}{136K}
\newcommand{\PeReachFp}{40K}
\newcommand{\PeReachFreeEE}{787}
\newcommand{\PeReachFreeEF}{709}
\newcommand{\PeReachFreeFF}{769}
\newcommand{\PeReachFreeFp}{619}
\newcommand{\PeReachPeakEE}{10,477}
\newcommand{\PeReachPeakEF}{10,555}
\newcommand{\PeReachPeakFF}{10,495}
\newcommand{\PeReachPeakFp}{10,645}
\newcommand{\PeRecurrentMiB}{2,048}
\newcommand{\PeRepHours}{14.2}
\newcommand{\PeRerunSame}{16}
\newcommand{\PeRestGiB}{7.2}
\newcommand{\PeRestMax}{7,412}
\newcommand{\PeRestMin}{7,406}
\newcommand{\PeRfLcppMtpDeep}{5.90}
\newcommand{\PeRfLcppMtpZero}{6.31}
\newcommand{\PeRfLcppNdDeep}{21.09}
\newcommand{\PeRfLcppNdZero}{24.98}
\newcommand{\PeRfNifDfDeep}{46.90}
\newcommand{\PeRfNifDfDeepMath}{92.25}
\newcommand{\PeRfNifDfZero}{56.66}
\newcommand{\PeRfNifDfZeroMath}{98.61}
\newcommand{\PeRfNifNdDeep}{41.23}
\newcommand{\PeRfNifNdZero}{46.99}
\newcommand{\PeRfPLcppMtpDeep}{1040.8}
\newcommand{\PeRfPLcppMtpDeepHi}{1074.5}
\newcommand{\PeRfPLcppMtpDeepLo}{1005.7}
\newcommand{\PeRfPLcppMtpZero}{1054.3}
\newcommand{\PeRfPLcppMtpZeroHi}{1079.9}
\newcommand{\PeRfPLcppMtpZeroLo}{1018.5}
\newcommand{\PeRfPLcppNdDeep}{109.4}
\newcommand{\PeRfPLcppNdDeepHi}{109.6}
\newcommand{\PeRfPLcppNdDeepLo}{109.2}
\newcommand{\PeRfPLcppNdZero}{91.4}
\newcommand{\PeRfPLcppNdZeroHi}{92.1}
\newcommand{\PeRfPLcppNdZeroLo}{90.4}
\newcommand{\PeRfPNifDfDeep}{40.6}
\newcommand{\PeRfPNifDfDeepHi}{42.8}
\newcommand{\PeRfPNifDfDeepLo}{36.7}
\newcommand{\PeRfPNifDfZero}{29.9}
\newcommand{\PeRfPNifDfZeroHi}{34.7}
\newcommand{\PeRfPNifDfZeroLo}{23.5}
\newcommand{\PeRfPNifNdDeep}{7.0}
\newcommand{\PeRfPNifNdDeepHi}{7.2}
\newcommand{\PeRfPNifNdDeepLo}{6.9}
\newcommand{\PeRfPNifNdZero}{2.7}
\newcommand{\PeRfPNifNdZeroHi}{5.6}
\newcommand{\PeRfPNifNdZeroLo}{1.2}
\newcommand{\PeRfWLcppMtpDeep}{67.33}
\newcommand{\PeRfWLcppMtpZero}{72.82}
\newcommand{\PeRfWLcppNdDeep}{44.17}
\newcommand{\PeRfWLcppNdZero}{47.81}
\newcommand{\PeRfWNifDfDeep}{65.95}
\newcommand{\PeRfWNifDfDeepMath}{85.92}
\newcommand{\PeRfWNifDfZero}{73.58}
\newcommand{\PeRfWNifDfZeroMath}{92.43}
\newcommand{\PeRfWNifNdDeep}{44.12}
\newcommand{\PeRfWNifNdZero}{48.25}
\newcommand{\PeRfXLcppMtpDeep}{11.41}
\newcommand{\PeRfXLcppMtpZero}{11.54}
\newcommand{\PeRfXLcppNdDeep}{2.09}
\newcommand{\PeRfXLcppNdZero}{1.91}
\newcommand{\PeRfXNifDfDeep}{1.41}
\newcommand{\PeRfXNifDfDeepProse}{1.69}
\newcommand{\PeRfXNifDfZero}{1.30}
\newcommand{\PeRfXNifDfZeroProse}{1.68}
\newcommand{\PeRfXNifNdDeep}{1.07}
\newcommand{\PeRfXNifNdZero}{1.03}
\newcommand{\PeRiseMax}{135}
\newcommand{\PeRiseMin}{94}
\newcommand{\PeRoomAll}{563}
\newcommand{\PeRoomNinetySix}{593}
\newcommand{\PeRoomSixtyFour}{593}
\newcommand{\PeRoomThirtyTwo}{593}
\newcommand{\PeShaOne}{cf8966ebc885}
\newcommand{\PeShaTwo}{2a72afda6bf0}
\newcommand{\PeSixDMax}{1.5}
\newcommand{\PeSixEngine}{035e227}
\newcommand{\PeSixFreedGqa}{1.62}
\newcommand{\PeSixFreedHyb}{0.72}
\newcommand{\PeSixHolm}{0.00211}
\newcommand{\PeSixI}{\ensuremath{+3.7}}
\newcommand{\PeSixIHi}{\ensuremath{+7.0}}
\newcommand{\PeSixILo}{\ensuremath{+0.3}}
\newcommand{\PeSixItems}{200}
\newcommand{\PeSixM}{\ensuremath{+3.7}}
\newcommand{\PeSixMHi}{\ensuremath{+10.2}}
\newcommand{\PeSixMHolm}{0.341}
\newcommand{\PeSixMLo}{\ensuremath{-2.8}}
\newcommand{\PeSixMargin}{10}
\newcommand{\PeSixStepGqa}{1.52}
\newcommand{\PeSixStepHyb}{1.03}
\newcommand{\PeSixStopAt}{190}
\newcommand{\PeSixVtMax}{66.0}
\newcommand{\PeSixVtMin}{8.0}
\newcommand{\PeSkipEstimate}{1.58}
\newcommand{\PeSkipLeft}{0.28}
\newcommand{\PeSlopeDev}{2.0}
\newcommand{\PeStageOneMin}{35.5}
\newcommand{\PeStageOneRss}{5.8}
\newcommand{\PeStageOneSteps}{13}
\newcommand{\PeStageThreeHours}{2.54}
\newcommand{\PeStageTwoHours}{21.72}
\newcommand{\PeStageTwoJobs}{5}
\newcommand{\PeTThreeErrTwentyEight}{8.2}
\newcommand{\PeTThreeErrZero}{13.2}
\newcommand{\PeTTwoKillN}{4}
\newcommand{\PeTTwoKillOf}{11}
\newcommand{\PeTTwoMaxErr}{27.6}
\newcommand{\PeTTwoMedErr}{9.2}
\newcommand{\PeTangoHOneHalf}{2.7}
\newcommand{\PeTangoHOneHi}{\ensuremath{+2.7}}
\newcommand{\PeTangoHOneLo}{\ensuremath{-2.7}}
\newcommand{\PeTangoHThreeHi}{\ensuremath{+0.9}}
\newcommand{\PeTangoHThreeLo}{\ensuremath{-8.1}}
\newcommand{\PeTangoHTwoHi}{\ensuremath{+5.3}}
\newcommand{\PeTangoHTwoLo}{\ensuremath{-11.6}}
\newcommand{\PeTangoHTwoSHi}{\ensuremath{+18.9}}
\newcommand{\PeTangoHTwoSLo}{\ensuremath{+4.8}}
\newcommand{\PeTangoKOneHi}{\ensuremath{+3.6}}
\newcommand{\PeTangoKOneLo}{\ensuremath{-3.6}}
\newcommand{\PeTangoKTwoHi}{\ensuremath{+18.9}}
\newcommand{\PeTangoKTwoLo}{\ensuremath{+4.8}}
\newcommand{\PeTitleMtpRate}{72}
\newcommand{\PeTitleRatioHi}{1.9}
\newcommand{\PeTitleRatioLo}{1.5}
\newcommand{\PeTokBytesPct}{3}
\newcommand{\PeTokGbExl}{6.63}
\newcommand{\PeTokGbIs}{7.94}
\newcommand{\PeTokGbIxxs}{6.84}
\newcommand{\PeTokPerCycleEE}{2.75}
\newcommand{\PeTokPerCycleFF}{2.77}
\newcommand{\PeTokPerCycleFp}{2.75}
\newcommand{\PeTorch}{2.10.0{+}cu128}
\newcommand{\PeTriton}{3.6.0}
\newcommand{\PeUseGapAbs}{18.8}
\newcommand{\PeUseGapDiff}{\ensuremath{-18.8}}
\newcommand{\PeUseGapHi}{\ensuremath{-4.2}}
\newcommand{\PeUseGapLo}{\ensuremath{-33.3}}
\newcommand{\PeUseGapN}{48}
\newcommand{\PeUseGapOtt}{15}
\newcommand{\PeUseGapTen}{6}
\newcommand{\PeVerifyAdoptHi}{1.92}
\newcommand{\PeVerifyAdoptLo}{1.86}
\newcommand{\PeVerifySOneHi}{1.52}
\newcommand{\PeVerifySOneLo}{1.47}
\newcommand{\PeVerifySTwoHi}{1.35}
\newcommand{\PeVerifySTwoLo}{1.30}
\newcommand{\PeVisionGiB}{0.86}
\newcommand{\PeWNineFitGain}{10.4}
\newcommand{\PeWNineGain}{7.4}
\newcommand{\PeWidthMax}{16.7}
\newcommand{\PeWidthMin}{13.5}



\title{\textbf{A 27-Billion-Parameter Model at Up to \PeTitleMtpRate{} Tokens per Second on a 2018 Gaming GPU: Turing Kernels Nearly Double exllamav3’s 2-Bit Decode Speed}}

\hypersetup{
  pdftitle={A 27-Billion-Parameter Model at Up to \PeTitleMtpRate{} Tokens per Second on a 2018 Gaming GPU: Turing Kernels Nearly Double exllamav3’s 2-Bit Decode Speed},
  pdfauthor={Matthew Schwartz},
  pdfsubject={Turing kernel patches improve original-weight Qwen3.8-27B decode speed on one stock RTX 2080 Ti, with bounded quality checks, 72K memory reach and tensor-core retrieval accuracy at 64K}
}

\author{Matthew Schwartz\singleauthorinfo{Independent researcher \enspace
\href{https://orcid.org/0009-0009-4171-7247}{ORCID 0009-0009-4171-7247}
\enspace \href{mailto:matthew.schwartz.research@gmail.com}{\texttt{matthew.schwartz.research@gmail.com}}}}
\date{November 2026}

\begin{document}
\maketitle

\begin{abstract}
\noindent
Fitting a large model in an older GPU's memory does not ensure fast decoding.
We adapted exllamav3's kernels to run the original Qwen3.8-27B weights at
2 bits per weight on a stock 11\,GiB RTX 2080 Ti, a 2018 consumer card.
Pre-registered rotated rounds compared the patched and stock engines on
three prompts. The patched engine decoded \PeFtWZero{} tokens/s undrafted,
\PeFtUndraftedRatio{} times stock, and \PeFtSTwoMtpZero{} with the model's
multi-token-prediction (MTP) head, \PeFtMtpRatio{} times stock; a second
session matched within 1\%. Most of the undrafted gain comes from an int8
path fitted to Turing's shared memory with unfused projections. Split-K
verification and tensor-core attention support drafting and filled contexts.
The patches changed no GSM8K outcome in the registered MTP check. On 500
GSM8K and 500 MMLU items, the stock engine's EXL3 weights scored
\PeQGsmExlNoAcc\% and \PeQMmluExlNoAcc\%; differences from llama.cpp's 2-bit
GGUFs were unresolved after adjustment. An 8-bit key-value cache held
\PeReachEE{} tokens, and the tensor-core kernel retained the reference
kernel's retrieval accuracy at 64K. These measurements show how kernels
adapted to an older card recover decode speed after low-bit weights make
the model fit. They cover our fastest measured configuration of the original
weights on one card and three prompts, not the fastest model on the card;
dated comparisons and detailed limits are in the appendix.
\end{abstract}


\section{Introduction}
\label{sec:intro}
A stock RTX 2080 Ti can hold the original weights and multi-token-prediction
(MTP) head of a 27-billion-parameter model at about 2 bits per weight.
The remaining constraint is execution: exllamav3 disables its int8 decode
path below 80\,KB of shared memory per block, while Turing provides 64\,KB;
its Triton attention kernel uses scalar arithmetic on this architecture
\citep{triton5066}. A model that fits can therefore leave substantial speed
unused.

We adapted the matrix and attention paths of exllamav3 to this card.
In pre-registered rotated rounds, the patched engine decoded
\PeFtWZero{} tokens/s undrafted and \PeFtSTwoMtpZero{} with MTP at an empty
context, \PeFtUndraftedRatio{} and \PeFtMtpRatio{} times the stock engine
(Figure~\ref{fig:speed}). A second session reproduced the rates within 1\%.
The main undrafted gain comes from an int8 path whose memory requests fit
the hardware, combined with unfused projections. Split-K verification
reduces the cost of checking drafts; tensor-core attention preserves speed
at filled contexts.

For a reader deciding whether an older consumer GPU remains useful, the
contribution is a measured configuration, credited kernel patches and a
reproducible explanation of their costs. We retain the original model
weights and test the faster build's GSM8K outcomes, descriptive quality
against llama.cpp's 2-bit files, filled-context memory reach and paired
retrieval accuracy. The int8 path, slice-height selection and split-K have
prior implementations; our attention kernel serves exllamav3's 8-bit cache
on Turing (Section~\ref{sec:related}). Methods and Results describe the
configuration and its measurements; Performance analysis explains the gains.
Appendix~\ref{app:null} preserves the comparison protocols, null results and
limits, followed by the registered prediction ledger and prompts.


\section{Methods}
\label{sec:setup}
A drafter proposes several next tokens; a verification pass evaluates them
as rows of one forward pass and retains accepted tokens. Here the drafter
is the model's MTP layer. We use arm for one configuration in a session and
anchor for a patched configuration repeated within that session. TabbyAPI
is the HTTP server used for the quality and long-context runs.


\subsection{Model and card}
\label{sec:model}
Qwen3.8-27B \citep{qwen2026qwen38card} has 64 layers: 16 full-attention
layers, each with 4 KV heads of 256 dimensions, and 48 Gated DeltaNet (GDN)
layers \citep{yang2025gated} that carry the sequence in a fixed-size recurrent
state.
One MTP layer \citep{gloeckle2024multitoken} drafts tokens for
self-speculative decoding. We used turboderp's EXL3 conversion at 2.0 bits per
weight, revision \q{SC\_2.00bpw\_H3} \citep{turboderp2026qwen38exl3sc}, whose
weight-file revisions and hashes are frozen in the artifact. Table~\ref{tab:files} lists every model file with
its size and average bits per weight: the EXL3 file's \PeFilesExlTrellisParams\,B
trellis-coded weights average \PeFilesExlTrellisBpw{} bits, and its embedding
table and vision encoder, both stored at 16 bits, bring the whole file to
\PeFilesExlBpw. EXL3 stores each tensor at a whole number of bits, chosen per
tensor: 1 to 4 in this file, most at 2. exllamav3 keeps a 16-bit embedding
table (\PeEmbedGiB\,GiB here) in host memory and copies to the card the one
row each token needs, as llama.cpp does with its input embeddings, and a
text-only run does not load the vision encoder (\PeVisionGiB\,GiB).
Everything else, the 64 layers, the output head and the MTP layer
(\PeOnCardGiB\,GiB of the file), sits on the card; that is what we mean by
the whole model on the card, or full offload.
Weights, the MTP layer and working buffers together occupy about
\PeRestGiB\,GiB; the remaining memory holds context.
% Generated by paper8/scripts/build_claims.py from testsuite/results/turing/model_files.json. Do not edit by hand.
\begin{table}[t]
\centering\small
\begin{tabular}{lL{3.4cm}rrrl}
\toprule
File & Source; used by & Parameters & GiB & Bits/weight & SHA-256\\
\midrule
EXL3 2.0\,bpw (two shards) & turboderp, revision SC\_2.00bpw\_H3; every exllamav3 arm; rafatxf's fork & 27.78\,B & 9.50 & 2.94 & \q{cf8966ebc885}\\
UD-IQ2\_XXS & unsloth; llama.cpp, NInferno & 26.90\,B & 6.77 & 2.16 & \q{e792d8fb3142}\\
UD-IQ2\_S & unsloth; llama.cpp with MTP; quality lanes & 26.90\,B & 7.80 & 2.49 & \q{7897d2c5a5ce}\\
MTP sidecar, full & unsloth; llama.cpp with MTP & 2.97\,B & 1.28 & 3.66 & \q{50d9ce5a6da3}\\
MTP sidecar, shared & ours, from unsloth's with the tables stripped; the quality lanes' MTP arm & 0.42\,B & 0.26 & 5.01 & \q{0dcbd3f1403a}\\
DFlash2 drafter & z-lab; llama.cpp and NInferno with DFlash2 & 1.92\,B & 1.06 & 4.71 & \q{1a25c56858e1}\\
\bottomrule
\end{tabular}
\caption{Model files. Parameters and average bits per weight are counted from each file's own headers (an EXL3 trellis tensor counts its $k\times n$ weights); sizes are on disk. The EXL3 file holds the MTP layer, which the GGUFs do not, and a 16-bit vision encoder that a text-only run does not load. The full MTP sidecar carries its own embedding and output tables; ours, the shared sidecar, strips them and uses the base model's. SHA-256 prefixes; every file but ours matched the Hub's LFS digest.}
\label{tab:files}
\end{table}

Every run used one RTX 2080 Ti (sm\_75, \PeCardMiB\,MiB, driver \PeDriver) in
a desktop at a 260\,W power limit, except the laptop-GPU measurements that
Section~\ref{sec:bytes} cites, which ran on an RTX 4080 Laptop GPU (Ada). By
stock we mean 11\,GiB of memory, unmodified firmware (no firmware change that
raises the card's tensor-core rate) and no raised power limit. All desktop GPU work ran as queued jobs, one at a
time, while other projects shared the desktop's CPU. The original measurements
ran between 30~October and 5~November~2026; the engine-version rerun ran on
12~November~2026 and the replacement comparison on 14~November.



\subsection{Turing kernel patches}
\label{sec:build}
The w13 build is exllamav3 1.6.5 \citep{turboderp2025exllamav3} with a chain
of patches and four dispatch settings. We wrote every patch except W13, which
is rafatxf's kernel; W13 is the only patch that borrows code. Build labels identify the last patch in the chain: w9 includes patches
through W9, w11b through W11b, and w13 through W13. The patches are listed
below with citations to prior work on the same or similar techniques:
\begin{itemize}\setlength{\itemsep}{1pt}
\item \textbf{Int8 matrix-vector products on 64\,KB parts (W7).} For
single-token decode, exllamav3 multiplies its trellis-coded weights with
activations quantized to int8 per slice. Upstream declines that path on any
GPU with less than 80\,KB of opt-in shared memory per block. W7 checks, per
kernel instance and device, whether the instance's largest request plus its
static shared memory fits, and runs the instances that fit. Running this
kernel class on 64\,KB Turing parts is prior art
\citep{vcruz3052026exl3turing,weicj2026turing,spiritbuun2026buun}.
\item \textbf{Slice height chosen by the host (W9).} The int8 kernel splits a
matrix's inner dimension into slices, multiplies each with its own thread
blocks and adds their partial sums at the end. W9 lets the host choose the
slice height, so that wide outputs such
as the 248K-column head no longer run as one slice with one block per
multiprocessor; partly prior art
\citep{ghazni2026rocm,yamz2026kyojin,spiritbuun2026buun}.
\item \textbf{Two-row verify passes on the int8 path (W10, W10b),} with a row
cap that fits 64\,KB; prior art \citep{weicj2026turing}.
\item \textbf{Tensor-core attention (W11, W11b):} our decode-attention kernel
for the 8-bit cache, called fd75d for flash decoding on sm\_75, Turing's
architecture (Section~\ref{sec:kernels}).
\item \textbf{A stepped row pointer in the int8 kernel (W12p):} the kernel's
row pointer steps by the row stride instead of being recomputed for each row,
which leaves its one- and two-row products bit-identical (its gain is in
Section~\ref{sec:verify}).
\item \textbf{Split-K matrix-vector products for verify passes of 3 to 16
rows (W13):} rafatxf's kernel \citep{rafatxf2026turing}, MIT-licensed, ported
unchanged. Mode~1 sends it the shapes rafatxf's own dispatch sends it;
mode~2, our change, sends it every shape.
\end{itemize}
The dispatch settings unfuse the attention layers' q, k and v projections
and the GDN layers' qkv and z projections through the upstream setting
\q{EXL3\_MGEMM\_N\_THRESHOLD=0}. Stock runs these bundles as one call on its
fp16 multi-matrix kernel. Only matrices 8,192 columns or wider, such as gate
and up, run as separate matrix-vector calls, using int8 where that path is
available. rafatxf's branch also lists unfusing among its Turing speedups.
We disable fd75d's live splits with \q{EXL3\_FD75D\_LIVE=0}; W11b's live
splits divide the filled context rather than the cache's allocated length.
The maximum int8 pass height (\q{EXL3\_INT8\_SQ\_MAX\_M}) and split-K mode
(\q{EXL3\_GEMV\_SK}) are both 1 undrafted and both 2 with MTP. Undrafted, the two split-K modes run the
same code, because only verify passes have three or more rows. Mode~2 became
the drafted configuration of record after a registered GSM8K check
(Section~\ref{sec:quality}).



\subsection{Decode measurements}
\label{sec:speedsetup}
The speed of record, F2, was pre-registered on 31~October at 13:42 EDT and
ran the same afternoon. Its rows span \PeFtFirst{} to \PeFtLast{} EDT. Its 14
arms (Table~\ref{tab:speed} in Appendix~\ref{app:engine-comparisons} lists them with their results) ran as fresh
processes in three rounds, in rotated order. Each process loaded the model with exllamav3's own loader and
generator (a 32K-token cache), warmed up on three prompts (code, math and
prose; Appendix~\ref{app:prompts}), then sampled up to 256 tokens from each
with exllamav3's default sampler (temperature 0.8, min-p 0.08, unseeded),
thinking off, either from the bare prompt or with 28,000 tokens of
deterministic filler ahead of it (28,050-token prompts); a run ends early if
the model emits its end-of-sequence token. Drafted arms used the MTP layer
with up to 6 draft tokens per step. The generator adapted the draft length
to the acceptance rate. Pooled speed is the total decoded tokens divided by the total decode wall
time across an arm's three prompts and three rounds
(\PeFtRows{} rows in all), so each prompt counts by the tokens it decoded.
Our arms used the 8-bit cache. The stock arms ran the official 1.6.5 wheel
(CUDA 12.8, torch 2.10.0); exllamav3 1.6.6, released on 31~October, leaves
every CUDA file on this model's decode path unchanged. The rafatxf arms ran rafatxf's Turing branch
\citep{rafatxf2026turing} at commit 15654d3 (exllamav3 1.5.3), built with
gcc~13 so that its fdq4 attention kernel compiled; the arms used its best
configuration, the 4-bit cache that fdq4 reads. An earlier session's rafatxf build had
silently fallen back to Triton attention; this session's log shows no
fallback. The registration required w13 to beat w11b undrafted by at least
3\% in every round (it did) and allowed the phrase ``w13 closes rafatxf's MTP
lead'' only if mode~1 won every round at an empty context; everything else is
reported as measured. Other projects raised the desktop's load average from
\PeFtLoadStart{} to \PeFtLoadEnd{} during the session; the card ran at 260\,W
and \PeFtTempLo--\PeFtTempHi\,$^\circ$C.

A confirmation session on 4~November repeated these anchors in three rotated
rounds. The session audit and other-engine protocols are retained in
Appendix~\ref{app:engine-comparisons}; speedups use comparisons within a session.


\subsection{Quality measurements}
\label{sec:qualitysetup}
On the same card we ran the EXL3 weights on the stock 1.6.5 wheel, undrafted
and with MTP, through TabbyAPI \citep{theroyallab2025tabbyapi}; the w13 build
with MTP in mode~2, and the w11b build undrafted; and unsloth's UD-IQ2\_S and
UD-IQ2\_XXS GGUFs of the same model \citep{unsloth2026qwen38gguf} through
llama.cpp's server \citep{gerganov2023llamacpp}. The undrafted GGUF lanes ran
llama.cpp at commit d63fbfe7. The drafted ones ran llama.cpp v0.4.0 (5266f24)
with our patch for a shared-embedding MTP sidecar: UD-IQ2\_S with unsloth's
MTP layer (up to 3 drafts) and UD-IQ2\_XXS with z-lab's DFlash2 block-diffusion
drafter \citep{chen2026dflash,zlab2026dflash2} (Q4\_K\_M, up to 7 drafts).
Both engines used an 8-bit cache and flash attention. Each lane answered 500
GSM8K items \citep{cobbe2021gsm8k} (a seeded sample of the 1,319-item test
set); the stock EXL3 and GGUF lanes also answered 500 MMLU items
\citep{hendrycks2021mmlu} (all 57 subjects). The patched checks used GSM8K
only. All lanes used
thinking on at medium effort, greedy decoding, at most 4,096 output tokens and
one request at a time, scored by one scorer on one response shape. The lanes
ran between 30~October and 1~November. Paired comparisons use exact McNemar
tests \citep{mcnemar1947}, Holm-adjusted \citep{holm1979} across the 14
same-card pairs of the batch. The w11b and w13 lanes had registered rules (no
more than 1.0 point below the configuration each replaced); the comparisons
with llama.cpp are descriptive. Time per answer is wall time per request.



\subsection{Context and retrieval measurements}
\label{sec:lcsetup}
Memory reach is the longest registered ladder step at which a fresh server
holds a filled prompt below \PeLimitMiB\,MiB, leaving \PeHeadroomMiB\,MiB
of the card unused. Each step fills all but \PeFillGapMin--\PeFillGapMax{}
tokens, decodes at most 64 tokens and samples memory once per second.
The cache formats and their byte costs are:
\paragraph{Cache formats.}
exllamav3 quantizes the cache in groups of 32 values: each group is rotated by
a 32-point Hadamard transform, scaled by its largest magnitude and stored at
$b$ bits with one fp16 scale, so each cached value costs $(b+\tfrac12)/8$ bytes. We
write a format as key bits,value bits: \fmt{8,8} (the 8-bit cache, the default
in our serving setup), \fmt{8,4} (the split cache), \fmt{4,4} (the 4-bit cache)
and \fmt{fp16}. With the MTP layer's cache, they cost \PeCacheMiBEE,
\PeCacheMiBEF, \PeCacheMiBFF{} and \PeCacheMiBFp\,MiB per 1,024 tokens (K).
During prefill, exllamav3 1.6.5 reads a quantized cache by staging up to 64K
tokens of it at a time as fp16, so the quantization acts on prefill as well as
decode.

\paragraph{Decode kernels.}
\label{sec:kernels}
exllamav3 decodes attention over its paged cache with a Triton kernel in the
flash-decoding pattern \citep{dao2023flashdecoding}: the keys are split across
thread blocks, each block keeps an online softmax \citep{milakov2018online},
and a combine kernel merges the partial results through their log-sum-exp.
For the 8-bit cache, w13 replaces the split phase with fd75d, a CUDA kernel
that keeps the splits, the combine and the precision (fp16 operands, fp32
accumulators and row sums, probabilities rounded to fp16 for the value
product) but computes its tiles on Turing's fp16 tensor cores. It reads the
packed 8-bit keys and values straight into registers, computes
$S^{\top}=KQ^{\top}$ with the grouped query rows in the small dimension, and
accumulates $O^{\top}\mathrel{+}=V^{\top}P^{\top}$. The other quantized
formats are read by Triton; the run records do not name the kernel that read
fp16, and prefill never uses fd75d. The
environment variable \q{EXL3\_FD75D=0} selects Triton for a whole server
process, so stage~2's 8-bit arms differ from its other arms in kernel as well
as in format; the accuracy arm separates the two.


Retrieval uses \PeNeedles{} seven-digit values for one key, inserted into
LongMemEval-S chat haystacks using RULER-style wording
\citep{wu2025longmemeval,hsieh2024ruler}. An answer is right only if it
contains every value. The same \PeHTwoN{} items run at 32K and 64K with
greedy decoding, thinking off and a 768-token answer allowance; none was
truncated. Calibration used separate items before the test set was built.
The direct kernel control compares Triton with the tensor-core kernel on
the same 8-bit cache at 64K. Paired accuracy differences have registered
90\% bootstrap intervals (10,000 item resamples): a lower bound above
$-10$ points is usable, and an interval within $\pm5$ is equivalent.
Single-arm accuracies use Wilson 95\% intervals \citep{wilson1927}.
Appendix~\ref{app:lc-methods} gives the full system envelope, calibration,
other cache comparisons, numerical checks and replication protocol.


\section{Results}
\label{sec:results}
\subsection{Decode throughput}
\label{sec:speed}
At an empty context, the patched build decoded \PeFtWZero{} tokens/s
undrafted against the stock wheel's \PeFtStockZero{}:
\PeFtUndraftedRatio{} times as fast, with per-round gains of
\PeFtUndraftedOverStockLo--\PeFtUndraftedOverStockHi\%.
With MTP the rates were \PeFtSTwoMtpZero{} against \PeFtStockMtpZero{},
\PeFtMtpRatio{} times as fast (\PeFtMtpOverStockLo--\PeFtMtpOverStockHi\%
per round). At 28K the patched build retained \PeFtWTwentyEight{} undrafted
and \PeFtSTwoMtpTwentyEight{} with MTP. The confirmation session reproduced its four rates within
\PeCfFtwoMin--\PeCfFtwoMax\% and supplied the missing stock 28K arms.
There, the pooled patched-to-stock ratios were \PeCfXStockNdDeep{}
undrafted and \PeCfXStockMtpDeep{} with MTP; the patched build was
faster in every round.

\begin{figure}[!tbp]\centering
\includegraphics[width=\linewidth]{\fig{fig_speed.pdf}}
\caption{Turing kernel patches nearly double undrafted decode speed and reach
\PeFtSTwoMtpZero{} tokens/s with MTP on a stock RTX 2080 Ti. Points pool
three prompts in three rotated rounds; lines span the round rates.
Filled points use empty context and open points 28K. Stock is exllamav3
1.6.5; the patched build uses the same weights and 8-bit cache. Its MTP
configuration uses all-shape split-K dispatch (Section~\ref{sec:build});
that setting does not change one-row undrafted code. F2 had no stock 28K arm.}
\label{fig:speed}
\end{figure}

Drafting gains depend on accepted proposals. At empty context, the patched
MTP build decoded \PeFtPpSTwoCode{}, \PeFtPpSTwoMath{} and
\PeFtPpSTwoProse{} tokens/s on code, math and prose; undrafted rates varied
by less than 0.5\%. Pooling weights prompts by the tokens actually decoded,
including early end-of-sequence finishes. Appendix~\ref{app:prompts} gives
per-prompt rates and lengths, and Appendix~\ref{app:engine-comparisons}
retains every measured build and the dated competitor comparisons.


\subsection{Quality and answer time}
\label{sec:quality}
The descriptive GGUF comparisons and answer-time measurements are in
Appendix~\ref{sec:answers} (Table~\ref{tab:quality}).

On GSM8K the faster builds stayed within their registered bound. w13 in
mode~2 with MTP answered every item as the stock wheel with MTP did: no discordant item, with the
same final answer on \PeQWSame{} of 500 and identical text on \PeQWIdent,
and against w11b with MTP, the registered reference, it scored the same
\PeQWModeTwoRefAcc\% with no discordant item and the same final answer on
\PeQWModeTwoRefSame{} of 500 (the registered bound was $-1.0$ point). Undrafted, the
w11b build scored
\PeQWElevenAcc\% (\PeQWElevenDiff{} points against stock; \PeQWElevenTen{}
against \PeQWElevenOtt, $p=\PeQWElevenP$), inside the same bound; a registered
bit-identity check found w13's one-row products identical to its predecessor's,
so undrafted w13 inherits that result. The MTP drafter changed no GSM8K
outcome on the stock wheel (\PeQMtpGsmTen{} against \PeQMtpGsmOtt) and
\PeQMtpMmluTen{} against \PeQMtpMmluOtt{} on MMLU. No MMLU lane ran on w13 in this 2080 Ti study.





\subsection{Context reach by cache format}
\label{sec:reach}
The 8-bit cache passed at 64K and 72K and failed at 80K, where it peaked at
\PePeakEEEighty\,MiB, \PeEEEightyOver\,MiB over the limit, so its reach is
\PeReachEE{} (Table~\ref{tab:reach}, Figure~\ref{fig:reach}a). The other
three formats passed every step of their ladders, so their reaches are lower
bounds: at least \PeReachEF{} for the split cache, \PeReachFF{} for the 4-bit
cache and \PeReachFp{} for fp16. Reach is a memory-fit result; paired retrieval accuracy was tested at 32K
and 64K, rather than at the 72K memory limit. The ladder and peaks appear in
Table~\ref{tab:reach}; failed projection predictions and longer-context
accuracy are retained in Appendix~\ref{app:cache-results}.
\begin{table}[!tbp]\centering\small\setlength{\tabcolsep}{5pt}
\begin{tabular}{@{}lrrlrl@{}}\toprule
Format & Bytes/token/layer & MiB per 1K tokens & Ladder & Reach & Peak at reach (free)\\\midrule
\fmt{8,8} & \PeBytesEE & \PeCacheMiBEE & 64K, 72K, 80K & \PeReachEE & \PeReachPeakEE\,MiB (\PeReachFreeEE)\\
\fmt{8,4} & \PeBytesEF & \PeCacheMiBEF & 64K, 88K, 96K & $\geq$\PeReachEF & \PeReachPeakEF\,MiB (\PeReachFreeEF)\\
\fmt{4,4} & \PeBytesFF & \PeCacheMiBFF & 64K, 96K, 128K, 136K & $\geq$\PeReachFF & \PeReachPeakFF\,MiB (\PeReachFreeFF)\\
\fmt{fp16} & \PeBytesFp & \PeCacheMiBFp & 32K, 36K, 40K & $\geq$\PeReachFp & \PeReachPeakFp\,MiB (\PeReachFreeFp)\\
\bottomrule\end{tabular}
\caption{Reach: the longest ladder step with a filled prompt whose peak stayed
at or under \PeLimitMiB\,MiB. Bytes per token in one attention layer; MiB per
1K tokens over the 16 attention layers and the MTP layer. Free: the card's
\PeCardMiB\,MiB minus the peak. The 8-bit cache failed at 80K
(\PePeakEEEighty\,MiB); the others passed their whole ladders.}
\label{tab:reach}
\end{table}



\subsection{Retrieval accuracy of tensor-core attention}
\label{sec:kernel}
On the same 8-bit cache at 64K, Triton and the tensor-core kernel each
answered \PeArmTritonK{} of \PeArmTritonN{} items correctly. Their paired
difference was \PeKOneDiff{} points (90\% interval
\PeKOneLo{} to \PeKOneHi{}), within the registered equivalence margin.
\PeKOneSame{} of \PeHTwoN{} answer texts were identical, bounding the
accuracy difference by \PeKOneBound{} points on these items.
The tensor-core kernel decoded at \PeArmEESixtyFourDecode{} tokens/s
against Triton's \PeArmTritonDecode{}, \PeDecodeRatio{} times as fast.
Appendix~\ref{app:kernel-control} retains the logit check and its 48K
trigger; these retrieval results apply to the tested 64K task.


\section{Performance analysis}
\label{sec:why}
Nsight Systems censuses explain where each build spends decode time.
They ran in separate sessions; their ratios describe kernel costs rather
than replacing the interleaved speed comparisons. The capture-to-session
spread is retained in Appendix~\ref{app:cost-model}.


\subsection{Kernel time breakdown}
\label{sec:census}
A registered census of the stock wheel, undrafted at an empty context, found
the card busy \PeCensusBusy\% of the decode time, with
\PeCensusIdleLo--\PeCensusIdleHi\,ms per token idle between kernels and about
\PeCensusKernels{} kernel launches per token (\PeCensusMs\,ms per token under
the profiler). The EXL3 matrix kernels took \PeCensusMatrix\% of kernel time:
the matrix-vector kernel \PeCensusGemv\%, the multi-matrix kernel that runs
fused projections \PeCensusMgemm\%, and the matrix-matrix kernel
\PeCensusGemm\%. Attention, the recurrent-state update and the norms share the
rest. At an empty context, matrix kernels offer more scope for reducing decode
time than host gaps or attention.

Table~\ref{tab:census} puts w13's drafted censuses beside the stock wheel's,
with the same prompts, profiler, report and grouping. With MTP at an empty
context, w13 took \PeCgWMs{} ms per token against \PeCgStockMtpMs{}
(\PeCgMsCut\% less, in captures a day apart): its matrix kernels took \PeCgMatrixCut\% less time and
attention \PeCgAttnCut\% less, while the small kernels took only
\PeCgSmallCut\% less. Matrix products still took \PeCgWMatrixShare\% of kernel
time, against \PeCgStockMtpMatrixShare\% on the stock wheel, so they remain the
largest component of kernel time after w13. At 28K, attention rose to \PeCgWDeepAttnShare\% of w13's
kernel time.

\begin{table}[!tbp]\centering\small\setlength{\tabcolsep}{4pt}
\begin{tabular}{@{}lrrrrr@{}}\toprule
Capture & ms/token & Matrix products & Attention & GDN state & Small kernels\\\midrule
stock 1.6.5, undrafted, empty & \PeCgStockNdMs & \PeCgStockNdMatrixUs{} (\PeCgStockNdMatrixShare\%) & \PeCgStockNdAttnUs{} (\PeCgStockNdAttnShare\%) & \PeCgStockNdGdnUs{} (\PeCgStockNdGdnShare\%) & \PeCgStockNdSmallUs{} (\PeCgStockNdSmallShare\%)\\
stock 1.6.5, MTP, empty & \PeCgStockMtpMs & \PeCgStockMtpMatrixUs{} (\PeCgStockMtpMatrixShare\%) & \PeCgStockMtpAttnUs{} (\PeCgStockMtpAttnShare\%) & \PeCgStockMtpGdnUs{} (\PeCgStockMtpGdnShare\%) & \PeCgStockMtpSmallUs{} (\PeCgStockMtpSmallShare\%)\\
w13 mode 2, MTP, empty & \PeCgWMs & \PeCgWMatrixUs{} (\PeCgWMatrixShare\%) & \PeCgWAttnUs{} (\PeCgWAttnShare\%) & \PeCgWGdnUs{} (\PeCgWGdnShare\%) & \PeCgWSmallUs{} (\PeCgWSmallShare\%)\\
w13 mode 2, MTP, 28K & \PeCgWDeepMs & \PeCgWDeepMatrixUs{} (\PeCgWDeepMatrixShare\%) & \PeCgWDeepAttnUs{} (\PeCgWDeepAttnShare\%) & \PeCgWDeepGdnUs{} (\PeCgWDeepGdnShare\%) & \PeCgWDeepSmallUs{} (\PeCgWDeepSmallShare\%)\\
\bottomrule\end{tabular}
\caption{Decode kernel time per decoded token by group, in \textmu s (share of
kernel time), under Nsight Systems: the three speed prompts, the 8-bit cache,
and MTP drafting up to 6 tokens where stated. ms/token is decode wall time per
token, the mean over the prompts. Matrix products are the EXL3 matrix
kernels, w13's split-K kernel included; attention includes fd75d. No stock
census ran at 28K. w13 is our full patch chain; mode 2 selects its
all-shape split-K dispatch (Section~\ref{sec:build}).}
\label{tab:census}
\end{table}



\subsection{Int8 decoding with 64\,KB of shared memory}
\label{sec:int8}
In the same census design, the build with W7's int8 path and unfused
projections took \PeIntEightAfter{} ms per token against \PeIntEightBefore{}
for the stock wheel, \PeIntEightChange\% less. Unfusing alone cut
\PeIntEightUnfused\% and the int8 path alone \PeIntEightAlone\%. They interact
because the fused projections run in the multi-matrix kernel, which the int8
path does not cover; unfused, they become matrix-vector products and move to
int8 too. In an interleaved session without the profiler, planned before W9
first ran, the host's choice of slice height (W9) then made the build
\PeWNineGain\% faster undrafted at an empty context, in every round, and
\PeWNineFitGain\% faster than the same build with that choice switched off.



\subsection{Verify passes and drafting}
\label{sec:verify}
Drafting turns most decode steps into verify passes of several rows, one per
drafted token plus one. On quantized weights, the verify pass's cost limits
what drafting gains \citep{zhang2025specmquant}. Composed from kernel benchmarks of each shape with the
census's mix of matrix shapes, a verify pass of 3 to 5 rows costs
\PeVerifyAdoptLo--\PeVerifyAdoptHi{} single-row passes' worth of matrix time
on w11b, whose int8 path stops at two rows and hands taller passes to the fp16
path. Split-K cuts that to \PeVerifySOneLo--\PeVerifySOneHi{} in mode~1 and
\PeVerifySTwoLo--\PeVerifySTwoHi{} in mode~2. Over w11b in F2, w13 in mode~1
was \PeFtSplitKMtpPct\% faster with MTP (mode~2, the configuration of record,
\PeFtModeTwoOverAdoptPct\%) but \PeFtSplitKUndraftedPct\%
undrafted at an empty context (\PeFtSplitKMtpTwentyEightPct\% and
\PeFtSplitKUndraftedTwentyEightPct\% at 28K), each in every round. The
undrafted gain is the stepped pointer, W12p.



\subsection{Attention kernel at long context}
\label{sec:depth}
At 28K, attention also accounts for substantial kernel time. In a census of w9, the build before fd75d,
Triton's split kernel took \PeDepthTritonUs\,\textmu s per token
(\PeDepthTritonCall\,\textmu s per call, one call per attention layer); with
fd75d the split phase took \PeDepthFdUs\,\textmu s (\PeDepthFdCall\,\textmu s
per call, \PeDepthCallRatio{} times faster), and the token fell from
\PeDepthBefore{} to \PeDepthAfter{} ms, \PeDepthChange\% less (two captures,
23 minutes apart; the token fell by less than the split phase saved because
its other kernels, the matrix kernels among them, ran longer in the second
capture). In registered
interleaved repeats without the profiler, undrafted, w11b decoded
\PeDepthTwoGain\% faster than w9 at 28K in every round (\PeDepthTwoW{}
against \PeDepthTwoFOne{} tokens per second), while the same w11b with fd75d
switched off gained only \PeDepthTwoOff\%, attributing the long-context gain
to fd75d.
In a kernel
benchmark at 28,050 tokens with one query token, Triton's split kernel took
\PeBenchDefault\,\textmu s at its default launch setting and \PeBenchBest\,\textmu s
at the best of \PeBenchConfigs{} settings, against \PeBenchFd\,\textmu s for
fd75d and a floor of \PeBenchFloor\,\textmu s to stream the layer's cache at
550\,GB/s. No Triton setting closes the gap, consistent with Triton emitting
scalar fused multiply-adds on Turing \citep{triton5066}; fd75d runs on the tensor
cores within \PeBenchFdOverFloor{} times the floor.



\section{Discussion}
\label{sec:discussion}
The gains come from matching kernel geometry and precision to the card:
unfused projections expose the int8 path, each instance stays within
64\,KB, split-K lowers verification cost, and tensor-core attention avoids
Triton's scalar path at filled contexts. This is an implementation and
measurement result for the original weights; it does not require a new
model or training.

For short prompts, MTP with all-shape split-K is the fastest configuration
we measured. The 8-bit cache with tensor-core attention combines
\PeReachEE{} memory reach with the tested 64K retrieval result.
For a single short answer over a long prompt, prefill remains important:
at 64K it ran near \PeArmEESixtyFourPrefill{} tokens/s and a request took
\PeArmEESixtyFourItemS{} seconds on average. Decode throughput matters most
for longer answers or follow-up turns that reuse the prompt.


\section{Related work}
\label{sec:related}
\textbf{Turing engines and kernels.} rafatxf's Turing branch of exllamav3
adds fdq4, tensor-core decode attention over the 4-bit cache, an fp16 split-K
matrix-vector kernel, unfused projections and prefill work, and reports 1.86
to 2.57 times over v1.5.3 on a 22\,GiB 2080 Ti \citep{rafatxf2026turing}. exllamav3's int8 path has run
on 64\,KB Turing parts in vcruz305's patch, weicj's branch and an EXL3 kernel
in a llama.cpp fork \citep{vcruz3052026exl3turing,weicj2026turing,spiritbuun2026buun};
host-chosen slice heights appear in AMD ports and kyojin
\citep{ghazni2026rocm,yamz2026kyojin}. On sm\_75, TurboMind decodes 8- and
4-bit caches on tensor cores \citep{lmdeploy2024turbomind}, a vLLM fork
decodes an FP8 cache for this model \citep{cookjohn2026vllmturing}, NInfer
reads a rotated int8 cache with int8 tensor cores \citep{ninfer2026titan},
and Brontir's NInferno, which builds on NInfer, runs this model's 2-bit GGUF on the stock
card \citep{brontir2026ninferno}.
Trellis-coded weights descend from QTIP \citep{tseng2024qtip};
ik\_llama.cpp decodes its trellis quantizations with int8 products on Turing
\citep{ikawrakow2025ikllama}, and Mirai~S runs a 2.4-bit trellis quantization
of this model on later cards \citep{trymirai2026mirais}. \textbf{Key and value precision.} Keys tolerate quantization worse than
values \citep{liu2024kivi,hooper2024kvquant}. QuaRot rotates caches with
Hadamard transforms \citep{ashkboos2024quarot}, as exllamav3 does per group of
32 values. Allocation studies give keys the extra bits
\citep{dong2024qaq,tao2024asymkv,zhang2024diffkv,hariri2026quantize}. In
KVTuner \citep{li2025kvtuner}, 8-bit keys with 4-bit values match an 8-bit
cache, and KV-COBRA \citep{ha2026kvcobra} finds retrieval the task that needs
key precision most. On this model, community llama.cpp tests found 4-bit
values nearly doubling tail divergence at 64K \citep{llamacpp23470} but
changing no short multiple-choice answers \citep{churnharness2026}, and FP8 or
NVFP4 caches keep its RULER scores near the full cache's to 64K and beyond on
data-center GPUs \citep{kozyrev2026gdn,tang2026jetqsparse}. Losses from
quantized caches grow with length \citep{li2024evaluating}, and quantization
moves the retrieval cliff shorter \citep{salammagari2026length}. In hybrids,
exact retrieval runs through the few attention layers
\citep{afendulev2026hybrid,bick2025gather}.

\textbf{Decode-kernel numerics.} Flash attention's deviation from standard
attention grows with sequence length \citep{golden2024flash}. FP8 tensor-core
sums cut long-context retrieval in vLLM until partial sums were promoted to
fp32 \citep{kubler2026fp8}, while Turing's fp16 tensor cores accumulate at
fp32 width with truncation \citep{fasi2021tensorcores}. Small kernel
differences change greedy outputs
\citep{yuan2025nondeterminism,du2026greedy,chen2026integer,he2025nondeterminism},
compressed models flip answers that an unchanged aggregate accuracy hides
\citep{dutta2024accuracy}, and evaluations under nominally deterministic settings vary in accuracy by up
to 15\% \citep{atil2024nondeterminism}. We compared fd75d with Triton on the
same cache by teacher-forced logits and by free-running accuracy.



\section{Artifact availability and reproducibility}
\label{sec:artifacts}
The prepared artifact includes frozen numeric outcomes, registrations,
analysis, figures and the credited patches. Its public identifier awaits
release. The model-free quick check, \q{python3 -B check\_claims.py},
rebuilds the claim macros and tables without a GPU, tokenizer or network;
\q{--figures} redraws the light figures. The replacement comparison uses a
hash-bound reviewed numeric projection, and the Bonsai tables replay nine
numeric lanes, five paired tests, Holm adjustment and fixed-seed bootstrap
intervals. These projections reproduce recorded scores, not response
rescoring or byte/token identity.

The full-item check, \q{scripts/repro\_check.sh}, also regenerates the
long-context prompts from committed manifests, using the pinned tokenizer,
LongMemEval-S data and dependencies. Full hardware replication needs the
revision-pinned EXL3 weights, GGUFs and drafters in Table~\ref{tab:files},
exllamav3 1.6.5, the patch overlay, TabbyAPI and an 11\,GiB sm\_75 card;
the artifact documents exact versions, commands and resource limits.
Weights and benchmark response text are linked or omitted rather than
redistributed. The patch set targets exllamav3 1.6.5 under MIT: W7--W12p
are ours; W13 is rafatxf's kernel from commit \q{15654d3}, ported unchanged
with its notice. The earlier 16K study's registration, outcomes and analysis
remain included.


\section{Conclusion}
\label{sec:conclusion}
Turing-fitted kernels run the original 2-bit Qwen3.8-27B weights on a stock
RTX 2080 Ti at \PeFtWZero{} tokens/s undrafted and \PeFtSTwoMtpZero{} with
MTP, \PeFtUndraftedRatio{} and \PeFtMtpRatio{} times stock exllamav3.
An int8 path within the card's shared-memory limit supplies most of the
undrafted gain; split-K verification and tensor-core attention support
MTP and filled contexts. The registered MTP check changed no GSM8K
outcome, while quality differences against llama.cpp's 2-bit GGUFs remain
unresolved after adjustment. An 8-bit cache holds \PeReachEE{} tokens and
the attention kernel retains reference retrieval accuracy at the tested
64K length. These results concern the fastest configuration we measured
for the original weights; they do not establish the fastest model on the
card. Appendix~\ref{app:null} preserves the null results, dated comparisons
and detailed limits.


\section*{AI assistance}
Claude Code (Anthropic) wrote the Turing patches W7 to W12p, including the
fd75d kernel, and ported rafatxf's split-K kernel (W13); wrote the
registrations, harnesses and analysis scripts; operated the runs under the
author's standing rules and registrations; and drafted the manuscript,
figures, claim checks and bibliography. fd75d was checked against an fp64
reference in the run that first timed it, W12p for bit identity, and the builds in
registered accuracy checks (Section~\ref{sec:quality}), with the test outputs
in the artifact; every number is recomputed from the archived
outputs by a script that refuses to build on any disagreement; and every
citation was checked against arXiv, Crossref, GitHub or the source page.
Codex (OpenAI) performed the registered Bonsai analysis and formatted its results;
in a separate review, independently recalculated the counts, rates and paired
tests and checked the timing boundaries and provenance. Codex also implemented the replacement speed harness, revised the manuscript
from the independently reviewed comparison, and prepared candidate reading
copies and the release package.
Matthew Schwartz directed the research and is responsible for its methods,
results, claims, citations, rights and released files.

\bibliographystyle{unsrtnat}
\bibliography{references}



\appendix
\section{Null results and limits}
\label{app:null}
This appendix retains incomplete and weaker comparisons, failed predictions,
replication outcomes and the detailed boundaries of the positive results.

\subsection{Engine comparisons at the measured versions}
\label{app:engine-comparisons}
\subsubsection{Protocols and session variability}
A registered audit of the \PeAuditSessions{} speed sessions
run with this harness found them \PeAuditChi{} times more variable than their
own round-to-round noise allows. The audit fits the logarithm of each
configuration's time per token in each session as a session effect plus a
configuration effect, weighting each cell by its round-to-round variance; the
factor is the weighted residual sum of squares (chi-squared) over its 33 degrees
of freedom, and 1 would mean the sessions agree within that noise. Three
sessions took \PeAuditSlowLo--\PeAuditSlowHi\% longer per token than F2, and
\PeAuditCloseWord{} others were within 2\% of it. The residuals' root mean square,
the cross-session floor, is \PeAuditFloorUndrafted\% undrafted and
\PeAuditFloorDrafted\% drafted. Speed claims are therefore comparisons
within one session, as performance studies elsewhere also advise
\citep{mytkowicz2009producing,hoefler2015scientific}.

A second session, registered on 4~November before it ran, repeated F2's
protocol and included the competing engines. Its conditions, each a fresh process in
three rotated rounds:
\begin{itemize}\setlength{\itemsep}{1pt}
\item the stock 1.6.5 wheel and w13 on 1.6.5, undrafted and with MTP, at an
empty context and 28K;
\item the 1.6.6 wheel and w13 rebuilt on 1.6.6, at an empty context;
\item rafatxf's branch with MTP;
\item NInferno \citep{brontir2026ninferno} at commit 1d3b43e3, undrafted on
both GGUFs;
\item llama.cpp v0.9.0 \citep{ggml2026llamacpp090}, undrafted on UD-IQ2\_XXS,
and with drafting: UD-IQ2\_XXS with DFlash2, and UD-IQ2\_S with unsloth's
full MTP sidecar;
\item ik\_llama.cpp \citep{ikawrakow2025ikllama} at 009dce19, undrafted and
with MTP;
\item buun-llama-cpp \citep{spiritbuun2026buun} at a8022d24, which loads the
EXL3 folder itself.
\end{itemize}
Every engine received the same prompts, filler, cap of 256 tokens and
nominal requested temperature and min-p settings. This did not establish
identical resolved sampling: defaults and the order of sampling operations
can differ by engine and version. The 4~November resolved presence penalty
is not established by our later audit. Timing spans the first decoded token
to the last; exllamav3 can stop at end of sequence, while the other engines'
servers decoded all 256. The exllamav3 builds ran in-process,
as in F2; the other engines ran as HTTP servers read by a streaming client,
whose timed span differed from each server's reported decode time by at most
\PeCfTimingMaxMs\,ms on all \PeCfTimingN{} rows with server timings
(\PeCfTimingCloseN{} within \PeCfTimingCloseMs\,ms). The analysis was fixed before any measurement was collected. There were four departures from the registration:
\begin{itemize}\setlength{\itemsep}{1pt}
\item FFmpeg, which NInferno links, had to be rebuilt without x86 assembly;
\item a memory stall caused by other programs on the desktop made its guard
stop the first attempt in its second round, and the queue restarted the
session from the beginning; the stopped attempt is not pooled;
\item seven competitor lanes failed before decoding, so a dated amendment,
written before any row was read, added a follow-up session of three rounds;
\item a registered 240\,W repeat never ran at 240\,W, since the power limit
could not be lowered without a password, so its rows are an unplanned repeat
at 260\,W.
\end{itemize}
The follow-up session tried each failed lane on a fit ladder with fixed rungs in
order:
\begin{itemize}\setlength{\itemsep}{1pt}
\item as registered;
\item a smaller prefill batch;
\item the smallest context that holds the 28K prompts;
\item the engine's own fit to free memory, reported as partial offload.
\end{itemize}
Same-session w13 anchors ran beside them. The session also used available ports, since the harness's port reuse had
caused two failures. In
both sessions the drafter for NInferno's DFlash2 was not on the machine, so
that arm did not run. It was downloaded on 5~November, at the revision
NInferno pins and matching the Hub's SHA-256, and two more dated amendments
ran it in sessions of three rotated rounds beside w13 anchors and NInferno
undrafted. The first session's lane could not start because NInferno serves a local
GGUF from its own text-only conversion and attaches no separate drafter to
it. The second session served NInferno's conversion of the same GGUF with the
drafter embedded, made beforehand by NInferno's own converter. With the
drafter's weights on the card, NInferno's memory plan cannot hold the
registered 32K context, and holds 29,696 tokens, the smallest that takes the
28K prompts, only on the last of its six plan rungs, a survival rung. That
rung keeps every weight and the whole context on the card but drops cached
prompt states and prefills 128 tokens at a time instead of 1,024; decode,
which we time, stays fully on the card. The lane therefore ran at NInferno's
automatic context when empty, as NInferno's published command does, and on
the survival rung at 28K.


\subsubsection{Registered comparison results}
The confirmation session's fork ran \PeCfRafVsFtPct\% faster than its F2
empty-context rate. The sessions therefore retain their own anchors.
Against rafatxf's branch with MTP, mode~2 was \PeFtOverRafZeroPct\% faster at
an empty context (\PeFtOverRafZeroLo\% to \PeFtOverRafZeroHi\% per round) and
\PeFtOverRafTwentyEightPct\% at 28K (\PeFtOverRafTwentyEightLo\% to
\PeFtOverRafTwentyEightHi\%), faster in every round at both. Mode~1 was
\PeFtSOneOverRafZeroPct\% slower at an empty context, in all three rounds, so
the registered phrase that w13 closes rafatxf's lead is not used. The two
builds ran different caches, each its best configuration: at an empty context
the cache is small and the comparison is mostly of matrix kernels; at 28K
rafatxf's 4-bit cache reads half the bytes of our 8-bit cache. Draft
acceptance was the same across the nine drafted arms (\PeFtAcceptLo--\PeFtAcceptHi{}
of drafted tokens, pooled), so their differences are kernel time. Some runs
ended at the end-of-sequence token before 256 tokens, most often on code
(Table~\ref{tab:f2prompts}), so the pooled rate counts each prompt by the
tokens it decoded. With the three prompts weighted equally instead, mode~2
led the fork by \PeFtOverRafZeroEqPct\% at an empty context and
\PeFtOverRafTwentyEightEqPct\% at 28K, and in the confirmation session by
\PeCfPRafZeroEq\% and \PeCfPRafDeepEq\%.

The confirmation session (Table~\ref{tab:confirm}) reproduced F2. w13's four
rates came within \PeCfFtwoMin\% to \PeCfFtwoMax\% of F2's, inside the
cross-session spread, so we retain F2 as the primary speed estimate. The stopped first attempt, not
pooled, agreed within \PeCfStopMax\%. In this session:
\begin{itemize}\setlength{\itemsep}{1pt}
\item exllamav3 1.6.6 changed no rate by more than \PeCfSixMax\%.
\item w13 led rafatxf's branch with MTP in every round, by \PeCfPRafZero\% at
an empty context and \PeCfPRafDeep\% at 28K.
\item Undrafted, w13 led NInferno in every round, on UD-IQ2\_XXS by
\PeCfPNifXxsZero\% and \PeCfPNifXxsDeep\% and on UD-IQ2\_S by
\PeCfPNifSZero\% and \PeCfPNifSDeep\%, at an empty context and 28K.
\item Undrafted, w13 ran \PeCfXLcppZero{} times as fast as llama.cpp v0.9.0 and
\PeCfXIkXxsZero{} times as fast as ik\_llama.cpp, which ran \PeCfIkVsLcpp\% slower
than llama.cpp on the same GGUF. ik\_llama.cpp turns CUDA graphs off below
Ampere (\texttt{ggml-cuda.cu} at the commit we built), and llama.cpp turned
them on for Turing \citep{llamacpp25749}: a likely cause, which we did not
test.
\end{itemize}
The follow-up session measured drafted competitors on the fit ladder:
\begin{itemize}\setlength{\itemsep}{1pt}
\item llama.cpp v0.9.0 with DFlash2 failed every rung, running out of
memory at load or on the first decode; its automatic fit cannot measure the
DFlash drafter's memory, a gap reported upstream the day before our session,
with a fix under review \citep{llamacpp31643}.
\item ik\_llama.cpp with MTP ran out of memory on its first request at all
three rungs that apply to it.
\item llama.cpp v0.9.0 with the MTP sidecar fit only by leaving layers in
host memory (the log does not record how many), and decoded at
\PeCfLcppMtpZero{} tokens per second.
\item buun-llama-cpp ran at full offload: \PeCfBuunNdZero{} tokens per second
undrafted (one round, before the time cap cut it) and \PeCfBuunMtpZero{} with
MTP, which w13 led by \PeCfXBuunMtpZero{} times in every round. Its
activations are int8; on 100 GSM8K items it scored \PeCfBuunAcc\% against the
stock EXL3 lane's \PeCfBuunRefAcc\% (\PeCfBuunTen{} items right only with
buun, \PeCfBuunOtt{} only with the stock lane; $p=\PeCfBuunP$).
\end{itemize}
In the 5~November DFlash2 session, NInferno with its drafter decoded \PeDfbNifDfZero{}
tokens per second at an empty context, \PeDfbNifDfGain{} times its undrafted
rate, and \PeDfbNifDfDeep{} at 28K. w13 with MTP decoded \PeDfbWMtpZero{} and
\PeDfbWMtpDeep, ahead by \PeDfbPNifDfZero\% and \PeDfbPNifDfDeep\% in every
round (\PeDfbPNifDfZeroLo\% to \PeDfbPNifDfZeroHi\% and \PeDfbPNifDfDeepLo\%
to \PeDfbPNifDfDeepHi\%), and its anchors came within \PeDfbFtwoMax\% of F2. The pooled
rate hides a split by prompt, which we did not register and report as
exploration (each prompt pooled over the three rounds). At an empty context
NInferno's DFlash2 kept pace on code (\PeDfbPpNifDfZeroCode{} against w13's
\PeDfbPpWMtpZeroCode) and math (\PeDfbPpNifDfZeroMath{} against
\PeDfbPpWMtpZeroMath), where it accepted \PeDfbAccZeroCode\% and
\PeDfbAccZeroMath\% of its drafted tokens. On prose it accepted
\PeDfbAccZeroProse\% and fell to \PeDfbPpNifDfZeroProse, below its own
undrafted \PeDfbPpNifXxsZeroProse, against w13's \PeDfbPpWMtpZeroProse.
NInferno drafts a fixed block of seven tokens; exllamav3 shortens its drafts
when acceptance falls. At 28K, w13 led on every prompt when pooled over rounds, though
NInferno's DFlash2 was faster on math in two of the three rounds. The lead
over NInferno's drafted lane is therefore a lead on this mix of prompts
under sampling, not on every kind of text.

\paragraph{Replacement comparison, 14 November.}
The 12~November rerun used NInferno 65e573fd and llama.cpp dd6b945b against
unchanged w13. We replaced its NInferno comparison: startup logs and the
pinned source reconstruct NInferno's default presence penalty of 1.5,
whereas exllamav3 used a neutral penalty. The 14~November replacement used
NInferno f8143b9e with presence penalty 0, frequency penalty 0 and repetition
penalty 1 asserted on every request; all 49 requests, including warmups and
the early diagnostic, resolved those values. The engine version also changed,
so differences between the earlier and replacement margins cannot be
attributed to the penalty. The reconstruction of the earlier default is not
an archived per-request receipt, and does not establish the 4~November setting.

The replacement kept the same card, frozen GGUFs (Hub revision 4ca72078,
Table~\ref{tab:files}), three prompts and three rotated rounds, with unchanged
w13 anchors. NInferno served its own UD-IQ2\_XXS conversion, loading the
DFlash2 drafter only when requested. Table~\ref{tab:refresh14} gives the four
eligible comparisons: w13 led undrafted by \PeFreshPNifNdZero\% empty and
\PeFreshPNifNdDeep\% at 28K, and with MTP against DFlash2 by
\PeFreshPNifDfZero\% and \PeFreshPNifDfDeep\%. Each pooled lead held in
every round. The empty-context undrafted margin is small and descriptive;
three prompts and three unpaired sampled rounds do not establish significance.
Neutral penalties do not make the samplers identical: NInferno's requested
top-k of 0 resolved to 64, with min-p applied after temperature within that
set; exllamav3 applied min-p before temperature. Seeds were not paired or
forwarded. The rates compare these configured stacks, without a fresh
quality-equivalence test or equal-output assumption.

Both code and prose favored w13 with drafting. At empty fill, code ran at
\PeFreshWNifDfZeroCode{} against \PeFreshNifDfZeroCode{} tokens/s and prose
at \PeFreshWNifDfZeroProse{} against \PeFreshNifDfZeroProse; at 28K the pairs
were \PeFreshWNifDfDeepCode{} against \PeFreshNifDfDeepCode{} and
\PeFreshWNifDfDeepProse{} against \PeFreshNifDfDeepProse. NInferno drafted
math faster at both fills: \PeFreshNifDfZeroMath{} against
\PeFreshWNifDfZeroMath{} tokens/s empty and \PeFreshNifDfDeepMath{} against
\PeFreshWNifDfDeepMath{} at 28K. These exploratory prompt splits limit the
pooled claim. NInferno's drafted context was 22,144 tokens when empty and
29,696 at 28K, against w13's 32,768; exllamav3 sometimes stopped before the
256-token cap, so equal requested caps did not mean equal generated work.

Only \PeFreshComplete{} of \PeFreshPlanned{} planned comparisons completed;
just the \PeFreshEligible{} NInferno comparisons qualified current
full-residency leads. The completed stock exllamav3 1.6.7 MTP comparison at
empty fill measured \PeFreshWStockMtpZero{} against \PeFreshStockMtpZero{}
tokens/s, but failed the registered all-module residency rule because its
embedding module was on CPU. This proves neither a compute-layer spill nor
an inability to fit. w13's existing CPU embedding policy was not freshly
audited by the same all-module rule. The other three stock pairs lacked
round-three cells and are not complete comparisons.

llama.cpp 61dbf492 decoded undrafted at \PeFreshLcppNdZero{} and
\PeFreshLcppNdDeep{} tokens/s and with its full MTP sidecar at
\PeFreshLcppMtpZero{} and \PeFreshLcppMtpDeep, at empty fill and 28K.
Its placement was unknown under the registered rule, so these are
measurements, not full-residency leads. DFlash2 had no successful measured
fit on the attempted ladder: all applicable rungs were tried and none
completed both warmups; four retained fit logs explicitly record CUDA
allocation failures. This is an unmeasured cell, not zero throughput or a
failure under every possible configuration. Full MTP failed the first three
rungs with allocation errors, then warmed up and measured at the last rung,
with context 29,696 and placement still unknown. Optional engines and newer
quantizations were not staged. This incomplete refresh does not revalidate
an all-engines or all-quantizations ranking. Earlier stock, fork, quality
and long-context results retain their original dates and scope.

% Generated from reviewed 14 November aggregate tables; descriptive ranges.
\begin{table}[!tbp]\centering\small\setlength{\tabcolsep}{5pt}
\begin{tabular}{@{}llrrll@{}}\toprule
Drafting & Fill & Patched EXL3 & NInferno & Excess (\%) & Round range (\%)\\\midrule
none / none & empty & \PeFreshWNifNdZero & \PeFreshNifNdZero & \PeFreshPNifNdZero & \PeFreshPNifNdZeroLo--\PeFreshPNifNdZeroHi\\
none / none & 28K & \PeFreshWNifNdDeep & \PeFreshNifNdDeep & \PeFreshPNifNdDeep & \PeFreshPNifNdDeepLo--\PeFreshPNifNdDeepHi\\
MTP / DFlash2 & empty & \PeFreshWNifDfZero & \PeFreshNifDfZero & \PeFreshPNifDfZero & \PeFreshPNifDfZeroLo--\PeFreshPNifDfZeroHi\\
MTP / DFlash2 & 28K & \PeFreshWNifDfDeep & \PeFreshNifDfDeep & \PeFreshPNifDfDeep & \PeFreshPNifDfDeepLo--\PeFreshPNifDfDeepHi\\
\bottomrule\end{tabular}
\caption{Replacement comparison, 14 November 2026: NInferno f8143b9e
against unchanged patched exllamav3, tokens per second pooled over three
prompts and three rotated rounds. Both use neutral penalties; top-k and
sampling order still differ. All four pooled leads hold in every round.
Ranges describe those rounds, not confidence intervals.}
\label{tab:refresh14}
\end{table}


\begin{table}[!tbp]\centering\small\setlength{\tabcolsep}{5pt}
\begin{tabular}{@{}llllrr@{}}\toprule
Build & Version & Drafting & Cache & Empty context & 28K tokens\\\midrule
stock & 1.6.5 wheel & none & \fmt{8,8} & \PeFtStockZero & --\\
w11b & 1.6.5 + chain to W11b & none & \fmt{8,8} & \PeFtAdoptZero & \PeFtAdoptTwentyEight\\
w13 & 1.6.5 + full chain & none & \fmt{8,8} & \PeFtWZero & \PeFtWTwentyEight\\\midrule
stock & 1.6.5 wheel & MTP & \fmt{8,8} & \PeFtStockMtpZero & --\\
w11b & 1.6.5 + chain to W11b & MTP & \fmt{8,8} & \PeFtAdoptMtpZero & \PeFtAdoptMtpTwentyEight\\
w13, mode 1 & 1.6.5 + full chain & MTP & \fmt{8,8} & \PeFtSOneMtpZero & \PeFtSOneMtpTwentyEight\\
w13, mode 2 & 1.6.5 + full chain & MTP & \fmt{8,8} & \PeFtSTwoMtpZero & \PeFtSTwoMtpTwentyEight\\
rafatxf & 1.5.3 + 15654d3 & MTP & \fmt{4,4} & \PeFtRafMtpZero & \PeFtRafMtpTwentyEight\\
\bottomrule\end{tabular}
\caption{Decode speed of record (F2): tokens per second, pooled over three
rotated rounds and three prompts, up to 256 tokens each, in one session on
31~October 2026. MTP: the model's MTP layer drafting up to 6 tokens per
step.
Stock had no 28K arm in this session. w11b and w13 are our patch chains
through the decode-attention and split-K patches, respectively; modes select
split-K dispatch (Section~\ref{sec:build}).}
\label{tab:speed}
\end{table}

\begin{table}[!tbp]\centering\small\setlength{\tabcolsep}{4pt}
\begin{tabular}{@{}llrrrr@{}}\toprule
Engine & Weights & \multicolumn{2}{c}{Undrafted} & \multicolumn{2}{c}{Drafted}\\
 & & empty & 28K & empty & 28K\\\midrule
\multicolumn{6}{@{}l}{\emph{Registered session, 4~November}}\\
exllamav3 1.6.5, stock & EXL3 2.0 & \PeCfStockNdZero & \PeCfStockNdDeep & \PeCfStockMtpZero & \PeCfStockMtpDeep\\
exllamav3 1.6.5 + w13 & EXL3 2.0 & \PeCfWNdZero & \PeCfWNdDeep & \PeCfWMtpZero & \PeCfWMtpDeep\\
exllamav3 1.6.6, stock & EXL3 2.0 & \PeCfSixNdZero & -- & \PeCfSixMtpZero & --\\
exllamav3 1.6.6 + w13 & EXL3 2.0 & \PeCfWSixNdZero & -- & \PeCfWSixMtpZero & --\\
rafatxf (4-bit cache) & EXL3 2.0 & -- & -- & \PeCfRafZero & \PeCfRafDeep\\
NInferno & UD-IQ2\_XXS & \PeCfNifXxsZero & \PeCfNifXxsDeep & -- & --\\
NInferno & UD-IQ2\_S & \PeCfNifSZero & \PeCfNifSDeep & -- & --\\
llama.cpp v0.9.0 & UD-IQ2\_XXS & \PeCfLcppZero & \PeCfLcppDeep & -- & --\\
ik\_llama.cpp & UD-IQ2\_XXS & \PeCfIkXxsZero & -- & -- & --\\
ik\_llama.cpp & UD-IQ2\_S & \PeCfIkSZero & -- & -- & --\\\midrule
\multicolumn{6}{@{}l}{\emph{Follow-up session with the fit ladder, 4~November}}\\
exllamav3 1.6.5 + w13 & EXL3 2.0 & \PeCfFitWNdZero & \PeCfFitWNdDeep & \PeCfFitWMtpZero & \PeCfFitWMtpDeep\\
llama.cpp v0.9.0 & UD-IQ2\_XXS & -- & -- & no fit & no fit\\
llama.cpp v0.9.0 & UD-IQ2\_S & -- & -- & \PeCfLcppMtpZero$^{\dagger}$ & \PeCfLcppMtpDeep$^{\dagger}$\\
ik\_llama.cpp & UD-IQ2\_S & -- & -- & no fit & --\\
buun-llama-cpp & EXL3 2.0 & \PeCfBuunNdZero$^{\ast}$ & -- & \PeCfBuunMtpZero & --\\\midrule
\multicolumn{6}{@{}l}{\emph{DFlash2 sessions, 5~November (Amendments 2 and 3)}}\\
exllamav3 1.6.5 + w13 & EXL3 2.0 & \PeDfaWNdZero & -- & \PeDfaWMtpZero & \PeDfaWMtpDeep\\
NInferno & UD-IQ2\_XXS & \PeDfaNifXxsZero & -- & no start & no start\\
exllamav3 1.6.5 + w13 & EXL3 2.0 & \PeDfbWNdZero & -- & \PeDfbWMtpZero & \PeDfbWMtpDeep\\
NInferno & UD-IQ2\_XXS & \PeDfbNifXxsZero & -- & \PeDfbNifDfZero$^{\ddagger}$ & \PeDfbNifDfDeep$^{\ddagger}$\\
\bottomrule\end{tabular}
\caption{The confirmation sessions: pooled decode tokens per second over three
rotated rounds, up to 256 tokens on each of the three prompts; each block is one
session with its own w13 anchors (our full patch chain; the DFlash2 block holds the two sessions of
Amendments 2 and 3), and comparisons stay within a block.
Drafted means the MTP head or sidecar, or DFlash2 for llama.cpp and NInferno
on UD-IQ2\_XXS. No fit: the lane failed every rung of the fit ladder. No
start: NInferno attaches no separate drafter to a local GGUF.
$^{\dagger}$ Fitted by the engine leaving layers in host memory, not a
full-offload speed. $^{\ast}$ One round, before the session's time cap cut
it. $^{\ddagger}$ On NInferno's own conversion of the same GGUF with the
drafter attached, at NInferno's automatic context of \PeDfbNifCtx{} tokens
when empty and \PeDfbNifDeepCtx{} at 28K (a survival rung).}
\label{tab:confirm}
\end{table}

\subsubsection{Version freshness at release preparation}
The comparisons above retain the versions and dates they measured.
The separate NInfer/Bonsai evaluation used NInfer
\q{f4ab369a06c728325b706c989d71be7dc0c05a1b}, tagged
\q{2026-11-12.3}. A newer release, \href{https://github.com/iamwavecut/ninfer-all/releases/tag/2026-11-16.2}{\q{2026-11-16.2}} at
\q{a3171940a3574cda4d4f3bc3000df7c0944e8ad8}, was published on
16~November 2026 Pacific time (17~November, 00:15 UTC) and was not measured
here. Its changes include a Turing position-split prefill shared-memory fix,
DSpine drafting, wider BF16 matrix-vector products and adaptive masked-draft
widths. NInferno had no formal release or tag at the 16~November check,
although its source had advanced beyond \q{f8143b9e}. The measured DFlash2
model payloads retained their recorded revisions and hashes. A comparison
against the newer NInfer release is follow-up work for the ternary decode
study, not evidence for a current engine ranking in this paper.

\subsection{Vendor reports and the ternary derivative}
\label{app:ternary}
The frozen EXL3 shard hashes begin \q{\PeShaOne} and \q{\PeShaTwo}; full
digests for these and the compared artifacts remain in the artifact inventory.
Public reports for this model on small cards were measured elsewhere, with
other harnesses. NInferno's own rows on this card give \PeNifUndrafted{}
tokens per second undrafted and \PeNifDflash{} with the DFlash2 drafter on
short chats, and \PeNifDeepUndrafted{} and \PeNifDeepDflash{} after a
\PeNifDeepTokens-token prompt, at 240\,W \citep{brontir2026ninferno}, as
listed when we surveyed it through 5~November; rows it posted on 8~November
give 48.29 undrafted on short chats and 43.30 over long prompts, at 260\,W
with a newer driver and its own timing; those rates are not directly comparable to ours. With
our prompts at 260\,W it decoded \PeCfNifXxsZero{} undrafted and, in a later
session, \PeDfbNifDfZero{} with DFlash2 (Table~\ref{tab:confirm}), and
\PeRfNifNdZero{} and \PeRfNifDfZero{} at its 12~November version
(the replaced presence-1.5 comparison). The replacement measured
\PeFreshNifNdZero{} and \PeFreshNifDfZero{} on 14~November.
For the derivative Ternary Bonsai~2 27B on this card, NInfer's vendor-reported October CLI runs list 58.5 tokens/s undrafted and 95.9 with drafting
\citep{ninfer2026smallcards}. These rates exceed our reported rates, but they
are not a matched P8 arm: prompts, harness and operating conditions differ. NInfer is a separate engine
from Brontir's NInferno \citep{brontir2026ninferno}. On later-generation cards with 10--12\,GiB, the original
weights decode at 45.1 tokens per second undrafted and 64.4 drafted (RTX 3080
10\,GiB, llama.cpp, a UD-IQ2\_XXS file its author lists at 9.0\,GB, larger
than the \PeFileGbIxxs\,GB file of that name we ran) \citep{sudoingx2026mtp}
and at 39.6 and 70.1 (RTX 4070 12\,GiB, Mirai~S at 2.4 bits per weight)
\citep{trymirai2026mirais}, at or below w13's rates here, though with other prompts
and harnesses; and the only other
original-weights report on a Turing card of 12\,GiB or less, QWENgine on an
8\,GiB Quadro RTX 4000 Mobile with some of its layers in host memory, gives
about 9 and 15 \citep{bensaad2026qwengine}.


\begin{table}[!tbp]\centering\footnotesize\setlength{\tabcolsep}{2.8pt}
\begin{tabular}{@{}>{\raggedright\arraybackslash}p{1.55in}*{4}{>{\raggedright\arraybackslash}p{1.08in}}@{}}\toprule
Counterpoint & None: GSM8K & None: MMLU & MTP-3: GSM8K & MTP-3: MMLU\\\midrule
Bonsai 2 27B/NInfer (ternary derivative; none / MTP-3) & $n=500$\newline 96.0\%\newline [93.9, 97.4]\%\newline 324.7 tokens\newline 6.02 s/answer\newline 9.56 correct/min & $n=500$\newline 88.4\%\newline [85.3, 90.9]\%\newline 679.3 tokens\newline 12.60 s/answer\newline 4.21 correct/min & $n=500$\newline 96.2\%\newline [94.1, 97.6]\%\newline 319.5 tokens\newline 2.53 s/answer\newline 22.84 correct/min & $n=500$\newline 87.8\%\newline [84.6, 90.4]\%\newline 684.6 tokens\newline 6.21 s/answer\newline 8.48 correct/min\\
\bottomrule\end{tabular}
\caption{One descriptive counterpoint row, outside the original-weight kernel
ablations: WaveCut/Ternary-Bonsai-2-27B-NInfer-v3 at revision \texttt{841f7812}
(artifact SHA-256 \texttt{cdc4810b\ldots c14a}), NInfer \texttt{f4ab369a}.
NInfer is a separate engine from Brontir's NInferno. All four cells have
500 requests, zero missing rows and zero transport errors; accuracy includes
wrong and length-capped answers, and brackets are Wilson 95\% intervals.
Tokens are mean native generated tokens (reasoning included), wall is mean
request time, and correct/min is 60 times correct answers divided by the sum
of all request wall seconds, excluding qualification and startup.
\emph{Normalized-shape footnote:} native split reasoning is scored through
P8's existing logical inline normalization and unchanged scorer; byte and
token identity to the original lanes is unverified. Length finishes were
0/4 for none and 0/3 for MTP-3 on GSM8K/MMLU; the MTP-3 MMLU count includes
one legacy-correct capped answer. The five registered paired accuracy
differences remain unresolved after Holm correction. The sentence's archived
w13 MTP control is \texttt{TURING-exl3-mtp-w13}; P8's mode-2 quality row uses
\texttt{TURING-exl3-mtp-w13s2}. No w13 MMLU control was registered.}
\label{tab:bonsai-counterpoint}
\end{table}

On our fixed 500-item GSM8K slice, Bonsai 2 27B/NInfer MTP-3 yielded 22.84 correct answers/min versus 13.03 for archived w13 MTP, at 96.2\% versus 97.2\% accuracy (paired difference -1.0 points, Holm-adjusted $p=0.680$); on our fixed 500-item MMLU slice, Bonsai 2 27B/NInfer MTP-3 yielded 8.48 correct answers/min versus 3.13 for stock EXL3 MTP, at 87.8\% versus 86.8\% accuracy (paired difference +1.0 points, Holm-adjusted $p=0.680$).

The vendor CLI rates above and the registered request-level table are different
measurements. The table's paired accuracy differences are unresolved; they
were not equivalence tests and do not establish indistinguishable quality.
\textbf{Long context on small cards.} rafatxf runs this model at 4.0 bits per
weight with a 4-bit cache on a modified 22\,GiB 2080 Ti, with single needles
found from 32K to 250K \citep{rafatxf2026turing}. On a stock 11\,GiB 2080 Ti,
NInfer running a ternary derivative allocates 200,704 tokens with a rotated 4-bit cache, or
143,360 with MTP \citep{ninfer2026smallcards}, and a registry entry recalls
one needle at 116,351 tokens \citep{sybil2026registry}. None measures accuracy
per cache format against a reference; our reach is a filled-context peak on
the original weights, and our accuracy is paired.


\subsection{Cross-card deployment replication}
\label{app:fleet}
% Fleet program parent: rife_f58bf36010d4b3dd. Result evidence:
% research/raw/2026-11-12-fleet-lanes/results-2026-11-13-7c2e/REPORT.md and
% research/raw/2026-11-12-fleet-lanes/results-2026-11-14-c4d8/REPORT.md.
\paragraph{A cross-card deployment replication, 13--14 November.}
Registered fleet tests ran the same EXL3 weights with w13 through TabbyAPI
HTTP on an RTX 4080 SUPER card and an RTX 4070 card
(Table~\ref{tab:fleet-replication}). All \PeFleetAssessedPredictions{} assessed
P8 predictions held: MTP exceeded undrafted speed at both fills on each card,
the 4080 SUPER was faster on the matched prompt mixture, GSM8K MTP scores
were within five percentage points of archived w13, and within-card
undrafted/MTP accuracy gaps were at most three points on either suite.
These are descriptive deployment replications, not equivalence tests or
speedup ratios against the 2080 Ti: HTTP scheduling/event coalescing,
stopping/batching and recurrent-cache settings differ from its direct client.
TabbyAPI did not forward request seeds; recorded seeds are nominal labels,
not controlled random streams. MMLU on w13 also extends the original
stock-build lane. The separate cross-card vote comparison remains unassessed.

% Generated by build_claims.py; committed fleet journals/verdict exports.
\begin{table}[!tbp]\centering\small\setlength{\tabcolsep}{4pt}
\begin{tabular}{@{}lrrrrrr@{}}\toprule
 & \multicolumn{4}{c}{HTTP decode tokens/s} & \multicolumn{2}{c}{Correct (none / MTP)}\\
Card & 0, none & 0, MTP & 28K, none & 28K, MTP & GSM8K & MMLU\\\midrule
RTX 4080 SUPER & \PeFleetCardANdZero & \PeFleetCardAMtpZero & \PeFleetCardANdDeep & \PeFleetCardAMtpDeep & \PeFleetCardAGsmNd/\PeFleetCardAGsmMtp & \PeFleetCardAMmluNd/\PeFleetCardAMmluMtp\\
RTX 4070 & \PeFleetCardBNdZero & \PeFleetCardBMtpZero & \PeFleetCardBNdDeep & \PeFleetCardBMtpDeep & \PeFleetCardBGsmNd/\PeFleetCardBGsmMtp & \PeFleetCardBMmluNd/\PeFleetCardBMmluMtp\\
\bottomrule\end{tabular}
\caption{Dated fleet deployment replication, reports of 13--14 November 2026.
Speed pools actual generated tokens over decode seconds for
\PeFleetSpeedRequestsPerCell{} requests per card/condition: three prompts in
three rotated rounds, warm-ups excluded. Quality is correct out of
\PeFleetQualityN{} items per suite/arm. Both cards serve the same EXL3 weights
with w13 over TabbyAPI HTTP; these are descriptive deployment results, not
speedup ratios against the 2080 Ti.}
\label{tab:fleet-replication}
\end{table}



\subsection{The earlier 16K cache study}
\paragraph{The earlier 16K study.}
\label{sec:p6}
The design builds on our earlier pre-registered study (unpublished; its
registration, data and analysis ship in this paper's artifact), which crossed
llama.cpp's q4\_0 cache with 2- and 4-bit GGUF weights of this model and of
the full-attention Qwen3-8B \citep{qwen2025qwen3} on \PeSixItems{} retrieval
items of about 16K tokens on a 12\,GiB laptop GPU. On the hybrid, the q4\_0
cache's cost in accuracy was larger at 2-bit than at 4-bit weights by
\PeSixI{} points (90\% interval \PeSixILo{} to \PeSixIHi), within that
study's equivalence margin of $\pm$\PeSixMargin{} points by two one-sided
tests (Holm-adjusted $p=\PeSixHolm$); the same 10-point margin is this
study's usability bound. Its memory measurements motivated the focus on context
length. At 16K the 4-bit cache freed \PeSixFreedHyb\,GiB on the hybrid, less
than the step from its 2-bit to its 3-bit weight file (\PeSixStepHyb\,GiB),
so reducing cache precision mainly permits longer contexts rather than
higher weight precision on this hybrid. Three of four task families were
answered fully in every configuration, which motivated difficulty
calibration in the present study.

\subsection{Long-context design and controls}
\label{app:lc-methods}
The long-context study comprised three stages, two decode-kernel checks
and a replication. Each was pre-registered before its first GPU step.
Stage~1 measured reach, stage~2 compared cache formats at 32K and 64K tokens,
and stage~3 ran the adopted format at 96K. A logit check and an accuracy arm
separated decode-kernel effects from cache-format effects. A replication
tested the large format difference on fresh items. Here, an arm is one cache
format at one length, run on a set of items in one TabbyAPI server session,
with the w13 build in its drafted configuration (Table~\ref{tab:envelope}).
Deviations were recorded as dated amendments before the steps they affected,
apart from one implementation note on where the logit check's analysis ran.

\paragraph{Reach.}
Stage~1 measured each format's reach: the longest length on a registered
ladder at which a server holding a filled prompt of that length peaked at most
\PeLimitMiB\,MiB, \PeHeadroomMiB\,MiB below the card's \PeCardMiB\,MiB. Each
step started a fresh server with context and cache sized to the step, sent one
prompt filling all but \PeFillGapMin--\PeFillGapMax{} tokens of it, decoded at
most 64 tokens and sampled the card's memory once per second. The ladders were
64K, 72K and 80K for \fmt{8,8}; 64K, 96K, 128K and 136K for \fmt{4,4}; 64K, 88K
and 96K for \fmt{8,4}; and 32K, 36K and 40K for \fmt{fp16}; a ladder stopped at
its first failure.

\paragraph{Items and scoring.}
The items come from the earlier study's generator, which follows RULER's
needle tasks
\citep{hsieh2024ruler} and builds each haystack from LongMemEval-S chat
sessions \citep{wu2025longmemeval}. A request is one user message with Qwen's
thinking mode off: a fixed instruction, the haystack with its needles, and a
question (Appendix~\ref{app:prompts}). The graded family is multi-value
retrieval: \PeNeedles{} needles give one key different seven-digit values,
spread through the haystack, and an answer is right only if it contains all
\PeNeedles. The generator keys its randomness by a registered seed and the
item identifier, so an item has the same key, values and question at every
length. Difficulty was fixed by calibration before any test item was built: on
separate items, the 8-bit cache at 64K answered \PeCalMvTwelve,
\PeCalMvTwentyFour{} and \PeCalMvFortyEight{} items at 12, 24 and 48 values,
and the 12-value level (\PeCalMvAcc\%), closest to the registered 65\%, was
chosen, after an amendment had replaced the registered multi-key and
variable-tracking families, which sat at ceiling and floor
(\PeCalMkSixteen, \PeCalMkFortyEight{} and \PeCalMkOneTwentyEight{} at 16, 48
and 128 keys; \PeCalVtTwo, \PeCalVtThree{} and \PeCalVtFour{} at 2, 3 and 4
links). The test set has \PeHTwoN{} items at each of 32K and 64K; stage~3 used
the \PeArmEFNinetySixN{} even-numbered items at 96K. Answers could use up to 768
tokens; the longest used \PeAnswerMax. A descriptive check set of 48 items
reuses the earlier study's single-needle, multi-key, multi-value and
variable-tracking questions.

\paragraph{Comparisons and statistics.}
Stage~2 ran \fmt{8,8} and \fmt{4,4} at 64K, and \fmt{fp16}, \fmt{8,8} and
\fmt{4,4} at 32K. Its comparisons are paired differences in the share of
items fully right: H2 (primary), \fmt{4,4} minus \fmt{8,8} at 64K; H1,
\fmt{8,8} minus \fmt{fp16} at 32K; and H3, \fmt{4,4} minus \fmt{fp16} at 32K.
Each has a 90\% percentile bootstrap interval over items (10,000 resamples). A
format is usable if the lower bound lies above $-10$ points, the earlier
study's margin, and equivalent if the whole interval lies within $\pm5$.
Exact McNemar tests give $p$-values, Holm-adjusted across H1 and H3 only, and
single accuracies carry 95\% Wilson intervals \citep{wilson1927}. If H2 was
not usable, a rescue rule ran \fmt{8,4} at 64K and tested it as H2s,
\fmt{8,4} minus \fmt{8,8}, with the same bound. Stage~3 ran the adopted
format at 96K. Because stage~2's 8-bit arms were read by fd75d and every other
arm by Triton, a logit check teacher-forced the gold answers on three items at
each of 28K, 48K and 64K through both kernels without drafting; divergence
from Triton at 48K or 64K triggered an accuracy arm, \fmt{8,8} at 64K read by
Triton, with K1 (primary), Triton minus fd75d on the same cache, and K2,
\fmt{8,4} minus \fmt{8,8} with Triton reading both.

\paragraph{Execution.}
Stage~1 took \PeStageOneMin{} GPU-minutes; stage~2, \PeStageTwoHours{} of its
budgeted 22 GPU-hours; stage~3, \PeStageThreeHours{} of 8; the logit check
\PeKlMinutes{} minutes; and the accuracy arm \PeArmHours{} hours. No arm was
interrupted, restarted or resumed; no item was excluded, timed out or
truncated; and the server's prompt-token count matched the generator's on
every item. The budget rule skipped one descriptive check-set arm.

\paragraph{Replication.}
\label{sec:replication}
A fifth registration, committed on 4~November before any of its arms ran,
tested the split cache's 64K lead again. Its primary pair of arms, R1, ran the
split cache and the 8-bit cache at 64K, both read by Triton, on \PeROneN{}
fresh items from a new seed. The lead counted as replicated if R1's 90\%
interval lay above 0. Three further arms measured the spread that a harmless
numerical change produces (N: the 8-bit arm rerun with its prefill attention
summed in a different order), the split cache against fp16 at 32K (R2) and
6-bit values at 64K (R3). K2's discordant items count as beyond the harmless
spread if the Wilson 95\% upper bound of N's discordance lies below K2's 17 of
96. Every arm closed with no item excluded, in \PeRepHours{} of the 15
budgeted GPU-hours; an fp16 check reproduced stage~2's fp16 arm on
\PeRTwoCheck{} of 16 items, so R2 reused that arm, as registered.

\begin{table}[!tbp]\centering\footnotesize\setlength{\tabcolsep}{4pt}
\renewcommand{\arraystretch}{1.06}
\begin{tabular}{@{}L{2.3cm}L{13.85cm}@{}}\toprule
Item & The long-context runs\\\midrule
Versions & exllamav3 \PeExlVersionShort{} with the w13 overlay on \q{PYTHONPATH} and its drafted settings (\q{EXL3\_MGEMM\_N\_THRESHOLD=0}, \q{EXL3\_FD75D\_LIVE=0}, \q{EXL3\_INT8\_SQ\_MAX\_M=2}, \q{EXL3\_GEMV\_SK=2}), journaled by stage~1, whose launcher stages~2--3 reuse; TabbyAPI commit \q{b1c2253}; torch \PeTorch, Triton \PeTriton, CUDA \PeCuda{} and Python \PePython. Items from the earlier study's frozen generator with registered seeds; committed manifests carry each item's SHA-256. Five registrations, changed only by dated amendments.\\
Processes, state & One TabbyAPI server per arm (\PeNStarts{} starts across stages~2--3 and the accuracy arm, none restarted), driven over HTTP by a separate harness process. The server owns the paged KV cache and the MTP layer's cache (context = cache size = the arm's length), the fp16 recurrent state and its checkpoints in a \PeRecurrentMiB\,MiB host-RAM cache, and a prefix cache, which the server reported unused at all \PeNRequests{} requests of stages~2--3 and the accuracy arm. Prefill in chunks of \PeChunk{} tokens; MTP drafting of up to \PeDraftTokens{} tokens per step, adjusted by the server's dynamic drafting.\\
Placement, memory & Weights (all but the 16-bit embedding table, which stays in host memory), caches and state on the card; everything but the cache took \PeRestMin--\PeRestMax\,MiB at 64K. The display and a streaming server held \PeBaselineMiB\,MiB before every session. Pass limit \PeLimitMiB\,MiB; card memory sampled at 1\,Hz. Server resident memory peaked at \PePeakRss\,GiB.\\
Decoding & Qwen's chat template with thinking off; temperature 0 and greedy sampling, logged at each of the \PeNRequests{} requests; at most 768 output tokens for graded items. Graded prompts left at least \PeRoomSixtyFour{} tokens of the context for the answer, and none stopped on length.\\
Failures & A time-out or server error excludes the item (listed, never scored wrong); an out-of-memory error stops the arm; 86\,$^\circ$C on three samples pauses the run. None occurred; the card peaked at \PeMaxTemp\,$^\circ$C.\\
Claims & Paired differences with 90\% percentile bootstrap intervals: usable if the lower bound lies above $-10$ points, equivalent if within $\pm$5. H2 and K1 primary; H1 and H3 Holm-adjusted; H2s under the rescue rule; reach as the longest passing step. Everything else is descriptive.\\
\bottomrule\end{tabular}
\caption{The long-context system. Values come from the run records
(journals, per-arm server configurations, server logs and the logit check's
run records), except the TabbyAPI commit and the weight revision, which come
from our serving-lane and artifact records.}
\label{tab:envelope}
\end{table}

\begin{figure}[!tbp]\centering
\includegraphics[width=\linewidth]{\fig{fig_reach.pdf}}
\caption{(a) Peak card memory with one filled prompt at each ladder step,
per cache format; the open cross marks the failed step. (b) Share of graded
items fully right per arm, with 95\% Wilson intervals, in the colours and
markers of (a); the open circle is the 8-bit cache read by Triton, the filled
\fmt{8,8} circles are read by fd75d. Arms at one length are paired on the same
items. Across lengths an item keeps its key, values and question but not its
haystack; only the 96K--64K comparison (Figure~\ref{fig:comparisons}) pairs
items across lengths.}
\label{fig:reach}
\end{figure}

\subsection{Cache-format comparisons}
\label{app:cache-results}
\label{sec:accuracy}
Table~\ref{tab:comparisons} and Figure~\ref{fig:comparisons} give the
registered comparisons, Table~\ref{tab:arms} the arms. At 32K the 8-bit cache
and fp16, read by an unrecorded kernel (Table~\ref{tab:arms}), answered the
same \PeArmEEThirtyTwoK{} of \PeArmEEThirtyTwoN{} items
correctly, with no discordant item and \PeHOneSameTexts{} of \PeHTwoN{} answers
identical in text: H1 is equivalent. The 4-bit cache scored \PeHThreeAbs{}
points below fp16 (\PeHThreeLo{} to \PeHThreeHi): usable, not shown
equivalent. At 64K the 4-bit cache answered \PeArmFFSixtyFourAcc\% and the
8-bit cache \PeArmEESixtyFourAcc\%; H2 is \PeHTwoDiff{} points (\PeHTwoLo{} to
\PeHTwoHi), so the 4-bit cache is not shown usable at 64K, and the data do not
resolve whether it costs anything there. Under the rescue rule the split cache
ran at 64K and answered \PeArmEFSixtyFourAcc\%: H2s is \PeHTwoSDiff{} points
(\PeHTwoSLo{} to \PeHTwoSHi; \PeHTwoSTen{} items right only with the split
cache, \PeHTwoSOtt{} only with the 8-bit cache), usable, and the rule adopted
it for stage~3. A rerun of the 8-bit arm on 16 items changed no outcome
(\PeNoiseK{} of \PeNoiseN), and the check set stayed at ceiling (single
needles, four keys or four values) or floor (variable tracking, 0 to
\PeCliffVtMax{} of 16) in all \PeCliffArmsWord{} arms that ran.

\begin{table}[!tbp]\centering\small\setlength{\tabcolsep}{3.4pt}
\begin{tabular}{@{}llrrrl@{}}\toprule
Comparison & Accuracy (a / b) & Difference & 90\% interval & Discordant & Verdict\\\midrule
H1: \fmt{8,8} $-$ \fmt{fp16}, 32K & \PeArmEEThirtyTwoAcc\% / \PeArmFpThirtyTwoAcc\% & \PeHOneDiff & [\PeHOneLo, \PeHOneHi] & \PeHOneTen{} / \PeHOneOtt & usable, equivalent\\
H3: \fmt{4,4} $-$ \fmt{fp16}, 32K & \PeArmFFThirtyTwoAcc\% / \PeArmFpThirtyTwoAcc\% & \PeHThreeDiff & [\PeHThreeLo, \PeHThreeHi] & \PeHThreeTen{} / \PeHThreeOtt & usable\\
H2: \fmt{4,4} $-$ \fmt{8,8}, 64K & \PeArmFFSixtyFourAcc\% / \PeArmEESixtyFourAcc\% & \PeHTwoDiff & [\PeHTwoLo, \PeHTwoHi] & \PeHTwoTen{} / \PeHTwoOtt & not shown usable\\
H2s: \fmt{8,4} $-$ \fmt{8,8}, 64K & \PeArmEFSixtyFourAcc\% / \PeArmEESixtyFourAcc\% & \PeHTwoSDiff & [\PeHTwoSLo, \PeHTwoSHi] & \PeHTwoSTen{} / \PeHTwoSOtt & usable; above $+5$\\
K1: \fmt{8,8} Triton $-$ fd75d, 64K & \PeArmTritonAcc\% / \PeArmEESixtyFourAcc\% & \PeKOneDiff & [\PeKOneLo, \PeKOneHi] & \PeKOneTen{} / \PeKOneOtt & usable, equivalent\\
K2: \fmt{8,4} $-$ \fmt{8,8}, Triton, 64K & \PeArmEFSixtyFourAcc\% / \PeArmTritonAcc\% & \PeKTwoDiff & [\PeKTwoLo, \PeKTwoHi] & \PeKTwoTen{} / \PeKTwoOtt & usable; above $+5$\\
R1: \fmt{8,4} $-$ \fmt{8,8}, 64K, fresh & \PeROneSplitAcc\% / \PeROneEightAcc\% & \PeROneDiff & [\PeROneLo, \PeROneHi] & \PeROneTen{} / \PeROneOtt & not replicated\\
N: \fmt{8,8} reordered $-$ \fmt{8,8}, 64K & \PeNzAccA\% / \PeNzAccB\% & \PeNzDiff & [\PeNzLo, \PeNzHi] & \PeNzTen{} / \PeNzOtt & equivalent\\
R2: \fmt{8,4} $-$ \fmt{fp16}, 32K & \PeRTwoAccA\% / \PeRTwoAccB\% & \PeRTwoDiff & [\PeRTwoLo, \PeRTwoHi] & \PeRTwoTen{} / \PeRTwoOtt & usable\\
R3: \fmt{8,6} $-$ \fmt{8,8}, 64K & \PeRThreeSixAcc\% / \PeRThreeEightAcc\% & \PeRThreeADiff & [\PeRThreeALo, \PeRThreeAHi] & \PeRThreeATen{} / \PeRThreeAOtt & descriptive\\
R3: \fmt{8,4} $-$ \fmt{8,6}, 64K & \PeRThreeFourAcc\% / \PeRThreeSixAcc\% & \PeRThreeBDiff & [\PeRThreeBLo, \PeRThreeBHi] & \PeRThreeBTen{} / \PeRThreeBOtt & descriptive\\
\bottomrule\end{tabular}
\caption{The registered paired comparisons on \PeHTwoN{} items each, in points
of items fully right; R1 uses fresh items. K2, R1, N and R3 have Triton
reading both arms; N reorders the prefill attention sums of the 8-bit arm. Discordant: items right only in arm a / only in arm b.
Exact McNemar $p$: H1 \PeHOneP, H3 \PeHThreeP{} (Holm \PeHThreeHolm), H2
\PeHTwoP, H2s \PeHTwoSP, K1 \PeKOneP, K2 \PeKTwoP, R1 \PeROneP, N \PeNzP, R2
\PeRTwoP, R3 \PeRThreeAP{} and \PeRThreeBP; all but H1's and H3's are
unadjusted. Usable: lower bound above $-10$; equivalent: interval
within $\pm5$ points. With few discordant items the percentile bootstrap is
narrow (H1's [0.0, 0.0]); Tango's score interval gives $\pm$\PeTangoHOneHalf{}
for H1, \PeTangoKOneLo{} to \PeTangoKOneHi{} for K1 and \PeTangoHTwoSLo{} to
\PeTangoHTwoSHi{} for H2s and K2; no usable or equivalent verdict changes
under it, though their lower bounds no longer clear $+5$.}
\label{tab:comparisons}
\end{table}

\begin{figure}[!tbp]\centering
\includegraphics[width=\linewidth]{\fig{fig_comparisons.pdf}}
\caption{Paired differences with 90\% bootstrap intervals. The hatched region
left of the dashed line marks differences below $-10$ points, where a lower
bound makes a format not usable; the light band marks the $\pm5$-point
equivalence region. The last row, the split cache at 96K against its own 64K
accuracy on \PeGapN{} items, is descriptive (Section~\ref{sec:ninetysix}).
The split cache's 64K lead (H2s) did not replicate: on fresh items, R1 gave
\PeROneDiff{} points (Section~\ref{sec:rep}), which this figure does not plot.}
\label{fig:comparisons}
\end{figure}

\begin{table}[!tbp]\centering\small\setlength{\tabcolsep}{4.5pt}
\begin{tabular}{@{}llrrrrr@{}}\toprule
Arm & Kernel & Items right (Wilson 95\%) & Decode tok/s & Prefill tok/s & s per request & Peak MiB\\\midrule
\fmt{fp16}, 32K & unrecorded & \PeArmFpThirtyTwoK{} (\PeArmFpThirtyTwoLo--\PeArmFpThirtyTwoHi\%) & \PeArmFpThirtyTwoDecode & \PeArmFpThirtyTwoPrefill & \PeArmFpThirtyTwoItemS & \PeArmFpThirtyTwoPeak\\
\fmt{8,8}, 32K & fd75d & \PeArmEEThirtyTwoK{} (\PeArmEEThirtyTwoLo--\PeArmEEThirtyTwoHi\%) & \PeArmEEThirtyTwoDecode & \PeArmEEThirtyTwoPrefill & \PeArmEEThirtyTwoItemS & \PeArmEEThirtyTwoPeak\\
\fmt{4,4}, 32K & Triton & \PeArmFFThirtyTwoK{} (\PeArmFFThirtyTwoLo--\PeArmFFThirtyTwoHi\%) & \PeArmFFThirtyTwoDecode & \PeArmFFThirtyTwoPrefill & \PeArmFFThirtyTwoItemS & \PeArmFFThirtyTwoPeak\\
\fmt{8,8}, 64K & fd75d & \PeArmEESixtyFourK{} (\PeArmEESixtyFourLo--\PeArmEESixtyFourHi\%) & \PeArmEESixtyFourDecode & \PeArmEESixtyFourPrefill & \PeArmEESixtyFourItemS & \PeArmEESixtyFourPeak\\
\fmt{8,8}, 64K & Triton & \PeArmTritonK{} (\PeArmTritonLo--\PeArmTritonHi\%) & \PeArmTritonDecode & \PeArmTritonPrefill & \PeArmTritonItemS & \PeArmTritonPeak\\
\fmt{8,4}, 64K & Triton & \PeArmEFSixtyFourK{} (\PeArmEFSixtyFourLo--\PeArmEFSixtyFourHi\%) & \PeArmEFSixtyFourDecode & \PeArmEFSixtyFourPrefill & \PeArmEFSixtyFourItemS & \PeArmEFSixtyFourPeak\\
\fmt{4,4}, 64K & Triton & \PeArmFFSixtyFourK{} (\PeArmFFSixtyFourLo--\PeArmFFSixtyFourHi\%) & \PeArmFFSixtyFourDecode & \PeArmFFSixtyFourPrefill & \PeArmFFSixtyFourItemS & \PeArmFFSixtyFourPeak\\
\fmt{8,4}, 96K & Triton & \PeArmEFNinetySixK{} of \PeArmEFNinetySixN{} (\PeArmEFNinetySixLo--\PeArmEFNinetySixHi\%) & \PeArmEFNinetySixDecode & \PeArmEFNinetySixPrefill & \PeArmEFNinetySixItemS & \PeArmEFNinetySixPeak\\
\bottomrule\end{tabular}
\caption{The graded long-context arms: items fully right of \PeHTwoN{} unless
stated, decode and prefill rates pooled over the arm's requests with MTP
drafting on, mean wall time per request, and the arm's peak card memory.
Speeds are descriptive and single-stream. fp16's decode rate reflects about
\PeCycleFpMs\,ms per draft-and-verify cycle, against \PeCycleEEMs\,ms for the
8-bit cache and \PeCycleFFMs\,ms for the 4-bit cache at 32K; drafting
accepted \PeAcceptFp\%, \PeAcceptEE\% and \PeAcceptFF\% of drafted tokens, so
drafting is not the cause. The run records do not name the attention path that
read the fp16 cache.}
\label{tab:arms}
\end{table}

\subsection{Effect of the decode kernel on accuracy}
\label{app:kernel-control}
The logit check found that the two kernels agreed on the most likely token at all \PeKlPositions{}
teacher-forced positions. Two Triton runs gave identical distributions, so the
registered divergence threshold reduced to $10^{-4}$ nats. The mean KL
divergence of fd75d from Triton was \PeKlTwentyEight{} nats at 28K,
\PeKlFortyEight{} at 48K and \PeKlSixtyFour{} at 64K, and the 48K result diverged by the registered rule
(the largest divergence at any one position was \PeKlMax{} nats), which
triggered the accuracy arm. There, Triton reading the 8-bit cache
answered \PeArmTritonK{} of \PeArmTritonN{} items correctly, the same count as
fd75d: K1 is \PeKOneDiff{} points (\PeKOneLo{} to \PeKOneHi), usable and
equivalent. \PeKOneSame{} of \PeHTwoN{} answer texts were identical between
the kernels, so $|K1|$ cannot exceed \PeKOneBound{} points on these items.
fd75d decoded at \PeArmEESixtyFourDecode{} tokens per second at 64K against
\PeArmTritonDecode{} for Triton on the same cache and items, \PeDecodeRatio{}
times as fast.

\subsection{Results at 96K}
\label{sec:ninetysix}
The split cache answered \PeArmEFNinetySixK{} of \PeArmEFNinetySixN{} items
correctly at 96K, \PeArmEFNinetySixAcc\%. On the same \PeGapN{} items it
answered \PeGapSixtyFourK{} at 64K, \PeGapAbs{} points more (\PeGapLo{} to
\PeGapHi{} for 96K minus 64K); the 8-bit cache answered \PeGapEEK{} of them at
64K, \PeGapEEAbs{} points above the 96K result (descriptive, across formats).
The registered usable context is measured against the format's own 32K
accuracy, which the replication's R2 arm later supplied: on the \PeUseGapN{}
shared items the split cache at 96K scored \PeUseGapAbs{} points below its
own 32K accuracy (90\% interval of the difference \PeUseGapLo{} to
\PeUseGapHi; \PeUseGapTen{} against \PeUseGapOtt{} discordant items), beyond
the registered 10 points, so its usable context ends below 96K. The 96K arm peaked at \PeArmEFNinetySixPeak\,MiB,
\PeNinetySixUnderLimit\,MiB under the pass limit, and each request took
\PeArmEFNinetySixItemS{} seconds on average, most of it prefill. 96K is a
length this card can hold; on this task, accuracy there fell more than 10
points below 64K.

\subsection{Replication of the split-cache result}
\label{sec:rep}
On the first item set, the split cache led the 8-bit cache at 64K by
\PeHTwoSAbs{} points, and with the kernel fixed (K2) the lead stayed:
\PeKTwoTen{} items right only with the split cache against \PeKTwoOtt{} (exact
McNemar $p=\PeKTwoP$). With the kernel held fixed, the registered design
attributes such a difference to the value bits plus numerical noise. On
\PeROneN{} fresh items, R1 gave \PeROneDiff{} points (\PeROneLo{} to
\PeROneHi; \PeROneTen{} against \PeROneOtt, $p=\PeROneP$): the lead did not
replicate, and the interval reaches below $-5$, so it is not equivalent to no
lead either. Greedy decoding turns small numerical differences into different
texts where the top two logits are close
\citep{yuan2025nondeterminism,du2026greedy}; changing the value bits changed
\PeKTwoDiffer{} of \PeHTwoN{} texts on the first set and \PeROneDiffer{} on the
fresh one. An imbalance in such flips on one item set can arise by chance; that set
alone cannot establish whether the imbalance generalizes. The 4-bit values showed no
accuracy advantage on fresh items. A cost of several points is not excluded (R1's interval
runs from \PeROneLo{} to \PeROneHi). A paired difference of this size between two long-context configurations
requires validation on fresh items.

The replication's other arms bound the contribution of numerical noise. A harmless change, rerunning the same 8-bit arm
with its prefill attention summed in another order (N), changed
\PeNzDisc{} of \PeNzN{} outcomes on the original items (N: \PeNzDiff{}
points, \PeNzLo{} to \PeNzHi), with a Wilson 95\% upper bound of
\PeNzWilsonHi\%, below K2's \PeNzKTwoDisc\%: the value bits change which
items succeed beyond numerical noise, but on fresh items not which way on
balance. At 32K the split cache answered as many original items as fp16
(\PeRTwoAccA\% each, \PeRTwoTen{} items each way; R2 \PeRTwoDiff{} points,
\PeRTwoLo{} to \PeRTwoHi), so it has no lead there, and the interval just
misses equivalence. On the original 64K items, 6-bit values (R3,
descriptive) scored \PeRThreeSixAcc\%, between 8-bit (\PeRThreeEightAcc\%)
and 4-bit values (\PeRThreeFourAcc\%). A graded pattern on one item set is what
perturbations that grow as the bits fall would produce if they flip the same
near-tie items; it is not evidence that fewer bits help. Of the registered
predictions, N's held (equivalent, at most 6 flips); R2's equivalence failed
narrowly, and R3's, that 6-bit values are equivalent to 8-bit, did not hold.


\subsection{Reach predictions and request costs}
Of stage~1's registered predictions, two
held (the 8-bit cache fit 64K, and the 4-bit cache reached at least 128K); the
third, that every peak would fall within $-150$ to $+250$\,MiB of a projection
from our earlier memory measurements, failed for the three quantized formats,
whose 64K peaks fell \PePOneMin--\PePOneMax\,MiB below their projections,
outside the band.
Beyond each ladder's first interval, the measured growth matched the cache
arithmetic within \PeSlopeDev\,MiB per 1K tokens. Over their 48 or 96
requests, the graded arms peaked \PeRiseMin--\PeRiseMax\,MiB above stage~1's
single-request peaks, and none exceeded \PeArmPeakMax\,MiB.


\subsection{Profiler scope and cost-model tests}
\label{app:cost-model}
Each census below ran the three speed prompts under Nsight
Systems, in its own session, so a comparison between censuses is also a
comparison between sessions. Against F2, the stock wheel's captures took
\PeProfStockPct\% longer per token undrafted and \PeProfStockMtpPct\% longer
with MTP, w13's drafted capture \PeProfWMtpPct\% less, and the undrafted
capture of w12p, the chain up to W12p, which runs w13's one-row kernels
unchanged when undrafted, \PeProfFinPct\% more. Ratios between censuses therefore overstate w13's
gain over the stock wheel (with MTP, \PeCgMsCut\% less time per token in the
censuses against \PeFtMtpCutPct\% in F2): we use censuses to
measure the distribution of time within a build and registered sessions to
measure speedups.

\subsection{A bandwidth model}
\label{sec:bytes}
If decoding were bound by memory bandwidth, a token would take the bytes it
reads (weights, recurrent state and cache) over the card's streaming rate. A
registered test compared this bandwidth model with a per-kernel cost model
at 598\,GB/s, 97\% of the card's rated 616\,GB/s. Bytes alone
underpredicted the undrafted w11b and w13 tokens of F2 by
\PeBytesErrMin--\PeBytesErrMax\%. The per-kernel model (the priced graph) sums each kernel class in a census.
It uses the larger of the time to read the class's bytes at the streaming rate
and its weight count times a per-class number of SM cycles per weight, then
adds launch costs and measured host idle time. Its constants were fitted on
declared configurations and tested on held-out ones. The test's claim that it would
come within $\pm12\%$ of every held-out
configuration was refuted (largest error \PeTTwoMaxErr\%, median
\PeTTwoMedErr\%; \PeTTwoKillN{} of \PeTTwoKillOf{} configurations missed by more
than 20\%), although it predicted time per token for F2's four undrafted w11b and w13 arms within
$+$\PePricedMin\% to $+$\PePricedMax\%.

Extended to drafting, by composing a draft-and-verify cycle from undrafted
kernel costs and the measured counts of drafted kernels, the priced graph
underpredicted w13 mode~2's time per
token by \PeTThreeErrZero\% at an empty context and \PeTThreeErrTwentyEight\% at
28K, outside the registered $\pm10\%$ at an empty context. Its matrix time
matches the drafted census; the missing time is in the non-matrix kernels,
including the extra recurrent-state work of the GDN layers that drafting
adds, and in the host gaps of each pass.

Matrix-call time beyond the weight-read time was measured
in a separately registered benchmark of w13's one-row int8 matrix
kernel on this card, whose registration and results ship in the artifact. The
bench ran without the profiler, warmed up before every timing and read each
weight tensor from memory, as decoding does. Because EXL3 stores each tensor
at a whole number of bits, it fitted each bit width separately, as a fixed
time per call plus a time per weight. A 2-bit call has a fixed time cost of
about \PeNineFixedTwo\,\textmu s, which CUDA-graph replay does not remove (an
empty kernel accounts for \PeNinePhNull\,\textmu s of it), plus a time per
weight \PeNinePerWeightTwoPct\% above the time to read the weight at
550\,GB/s, 89\% of the rated rate and a little below the
\PeNineFourBitGbs\,GB/s at which the bench's 4-bit calls read their weights.
Applied to the calls of the undrafted w12p census, those fits put the fixed
parts at \PeNineShare\% of the matrix time above the read (\PeNineShareRated\%
at the rated rate). A 1-bit call takes
\PeNineOneTwoLo--\PeNineOneTwoHi\% of a 2-bit call's time on six of the
model's seven matrix shapes (\PeNineOneTwoKv\% on the smallest): a 1-bit
weight takes \PeNineOneTwoPerWeight\% of a 2-bit weight's time,
\PeNinePerWeightOneX{} times as long as reading it, so fewer bits would barely
speed this card's matrix calls. rafatxf reported the same for whole tokens on
Turing, from one comparison at 3.5 and 4.0 bits per weight
\citep{rafatxf2026turing}. A registered repeat of the undrafted census on an
RTX 4080 Laptop GPU (Ada), with the stock engine and unfused projections,
found 1-bit calls quicker there: a 1-bit call took
\PeAdaOneFasterLo--\PeAdaOneFasterHi\% less time than a 2-bit call in the same
role at 80\,W and \PeAdaBoostOneFasterLo--\PeAdaBoostOneFasterHi\% less at
175\,W (Appendix~\ref{sec:predictions} lists its predictions). These
measurements are reported here, not analysed further.

\subsection{Time per answer}
\label{sec:answers}
Table~\ref{tab:quality} gives the lanes. On GSM8K every lane scored
\PeQGsmMin--\PeQGsmMax\%; against EXL3 undrafted, UD-IQ2\_S scored
\PeQPGsmIsAbs{} point lower (\PeQPGsmIsTen{} items right only with EXL3,
\PeQPGsmIsOtt{} only with the GGUF) and UD-IQ2\_XXS \PeQPGsmIxxsAbs{} lower
(\PeQPGsmIxxsTen{} against \PeQPGsmIxxsOtt). On MMLU, UD-IQ2\_S was
\PeQPMmluIsAbs{} point below EXL3 (\PeQPMmluIsTen{} against \PeQPMmluIsOtt,
$p=\PeQPMmluIsP$) and UD-IQ2\_XXS \PeQPMmluIxxsAbs{} points below
(\PeQPMmluIxxsTen{} against \PeQPMmluIxxsOtt, $p=\PeQPMmluIxxsP$, Holm
$p=\PeQPMmluIxxsHolm$); with drafting, IQ2\_XXS with DFlash2 was
\PeQPMmluDfAbs{} points below EXL3 with MTP (\PeQPMmluDfTen{} against
\PeQPMmluDfOtt, Holm $p=\PeQPMmluDfHolm$). The IQ2\_XXS gap is suggestive, not
significant after adjustment.

With drafting, w13 answered a GSM8K item in \PeQGsmExlWWall\,s on average,
against \PeQGsmIsMtpWall\,s for IQ2\_S with MTP and \PeQGsmIxxsDfWall\,s for
IQ2\_XXS with DFlash2: \PeQRatioIsMtp{} and \PeQRatioIxxsDf{} times faster per
answer. Section~\ref{sec:answers} decomposes the gap. The lanes ran in
different sessions, so these times are descriptive.

\begin{table}[!tbp]\centering\small\setlength{\tabcolsep}{3.6pt}
\begin{tabular}{@{}lllrrrrr@{}}\toprule
Lane & Engine & Drafter & GSM8K & MMLU & Tokens & tok/s & s/answer\\\midrule
EXL3 2.0 bpw & exllamav3 1.6.5 & none & \PeQGsmExlNoAcc & \PeQMmluExlNoAcc & \PeQGsmExlNoTok & \PeQGsmExlNoTps & \PeQGsmExlNoWall\\
EXL3 2.0 bpw & 1.6.5 + w11b & none & \PeQWElevenAcc & -- & \PeQWElevenTok & \PeQWElevenTps & \PeQWElevenWall\\
UD-IQ2\_S & llama.cpp d63fbfe7 & none & \PeQGsmIsAcc & \PeQMmluIsAcc & \PeQGsmIsTok & \PeQGsmIsTps & \PeQGsmIsWall\\
UD-IQ2\_XXS & llama.cpp d63fbfe7 & none & \PeQGsmIxxsAcc & \PeQMmluIxxsAcc & \PeQGsmIxxsTok & \PeQGsmIxxsTps & \PeQGsmIxxsWall\\\midrule
EXL3 2.0 bpw & exllamav3 1.6.5 & MTP & \PeQGsmExlMtpAcc & \PeQMmluExlMtpAcc & \PeQGsmExlMtpTok & \PeQGsmExlMtpTps & \PeQGsmExlMtpWall\\
EXL3 2.0 bpw & 1.6.5 + w11b & MTP & \PeQWElevenMtpAcc & -- & \PeQWElevenMtpTok & \PeQWElevenMtpTps & \PeQWElevenMtpWall\\
EXL3 2.0 bpw & 1.6.5 + w13, mode 2 & MTP & \PeQGsmExlWAcc & -- & \PeQGsmExlWTok & \PeQGsmExlWTps & \PeQGsmExlWWall\\
UD-IQ2\_S & llama.cpp v0.4.0 & MTP & \PeQGsmIsMtpAcc & \PeQMmluIsMtpAcc & \PeQGsmIsMtpTok & \PeQGsmIsMtpTps & \PeQGsmIsMtpWall\\
UD-IQ2\_XXS & llama.cpp v0.4.0 & DFlash2 & \PeQGsmIxxsDfAcc & \PeQMmluIxxsDfAcc & \PeQGsmIxxsDfTok & \PeQGsmIxxsDfTps & \PeQGsmIxxsDfWall\\
\bottomrule\end{tabular}
\caption{Same-card lanes: accuracy (\%) on 500 GSM8K and 500 MMLU items;
on GSM8K, mean tokens per answer (thinking included), median decode rate in
tokens per second and mean wall time per answer in seconds. One request at a time, greedy, thinking on. Each
lane ran in its own session (a few items of the stock lanes came from earlier
attempts); no MMLU lane in this 2080 Ti study ran on w11b or w13. w11b and w13 are our patch
chains through decode attention and split-K, respectively (Section~\ref{sec:build}).}
\label{tab:quality}
\end{table}

GSM8K prompts are short, so time per answer is close to the token count divided by the decode
rate. Against IQ2\_S with MTP, w13 generated \PeQTokRatioIsMtp{} times fewer
tokens per answer (means) and decoded \PeQTpsRatioIsMtp{} times faster
(medians), which
together account for the measured \PeQRatioIsMtp{} times shorter answer time;
against IQ2\_XXS with DFlash2 the factors are \PeQTokRatioIxxsDf{} and
\PeQTpsRatioIxxsDf, against \PeQRatioIxxsDf{} measured. Undrafted, the stock
wheel and llama.cpp decoded at about the same rate (\PeQGsmExlNoTps{} against
\PeQGsmIsTps{} and \PeQGsmIxxsTps{} tokens per second), so the stock wheel's
shorter answer time (\PeQGsmExlNoWall{} against \PeQGsmIsWall{} and
\PeQGsmIxxsWall{} s) comes from reasoning length: on the same items the 2-bit
GGUFs thought \PeQLongerMin--\PeQLongerMax\% longer than the 2-bit EXL3
weights, as low-bit quantization lengthens reasoning in other models
\citep{liu2025quantization}. Differences between weight formats therefore account for part of the
per-answer gap against llama.cpp, independently of engine speed. The drafted llama.cpp lanes ran v0.4.0; v0.9.0 adds
CUDA graphs for MTP drafts and device copies for DFlash
\citep{ggml2026llamacpp090}, so the per-answer ratio is against the builds we
measured. On v0.9.0 at the confirmation's 32K context, the DFlash2 lane did
not fit on the card at full offload, and the MTP lane, with unsloth's full
1.28\,GiB sidecar, fit only with layers in host memory
(Appendix~\ref{app:engine-comparisons}). The per-answer lanes fit at smaller contexts,
8,192 tokens for DFlash2 and 16,384 for MTP, and the MTP lane used our
0.26\,GiB shared-embedding sidecar.


\paragraph{Empty-context attention control.}
At an empty context, where attention is small, w11b was
\PeDepthTwoZeroGain\% faster than w9, in two of the three rounds, outside the
registered $\pm3\%$. \subsection{Contributions and prior work}
\label{sec:new}
\textbf{Measurement contribution.} To our knowledge (public web, GitHub, Hugging Face
and forks searched through 5~November 2026; Reddit, X, Zhihu and Bilibili only
through search-engine snippets and mirrors; Discord and video not at all),
this is the first published EXL3 measurement of the unmodified Qwen3.8-27B
weights on a single stock 11\,GiB RTX 2080 Ti, and the first with every
layer, the output head and the model's own MTP layer on one Turing card of
12\,GiB or less (Section~\ref{sec:model}: only the embedding table, read one
row per token, stays in host memory).
NInferno measured the unmodified weights on this card first, as a 2-bit GGUF
undrafted and with DFlash2 \citep{brontir2026ninferno}. QWENgine ran the MTP
head on an 8\,GiB Turing card with some of its layers in host memory
\citep{bensaad2026qwengine}, and a fork of exllamav3 split the model at
4.0 bits per weight across a 6\,GiB RTX 2060, which held the MTP layer, and a
16\,GiB Tesla P100 \citep{trisert2026pascal}. The dated sessions establish the comparisons in Section~\ref{sec:speed} and Appendix~\ref{app:engine-comparisons}, within their measured versions and conditions. They cover only the tested configurations and
do not rank all engines or quantizations. In the replacement of 14~November, w13 led
NInferno f8143b9e undrafted by \PeFreshPNifNdZero\% empty and
\PeFreshPNifNdDeep\% at 28K, and with MTP against DFlash2 by
\PeFreshPNifDfZero\% and \PeFreshPNifDfDeep\%, pooled in every round.
Code and prose favored w13; math favored NInferno at both fills.
Only \PeFreshComplete{} of \PeFreshPlanned{} comparisons completed and
\PeFreshEligible{} qualified current full-residency leads. Missing and
placement-unknown cells cannot support a wider ranking. Older measurements
of other engines remain dated observations. The ternary derivative and the
later-generation cards in Appendix~\ref{app:ternary} provide further context,
not controlled cross-model or cross-card comparisons. The Tesla T4 has the
same 64\,KB of shared memory as this card; we did not measure one.

\textbf{fd75d.} To our knowledge it is the first tensor-core decode-attention
kernel for exllamav3's 8-bit quantized cache on Turing. Its techniques have prior art: LMDeploy's TurboMind decodes 8- and 4-bit caches on Turing's tensor
cores with the same $S^{\top}=KQ^{\top}$ orientation
\citep{lmdeploy2024turbomind}. rafatxf's fdq4 reads exllamav3's 4-bit cache
straight into tensor-core fragments on Turing \citep{rafatxf2026turing}.
kyojin's \q{qdec} is a drop-in split kernel for exllamav3's 2- to 8-bit caches
on later GPUs \citep{yamz2026kyojin}. llama.cpp uses the same orientation
for fp16 caches \citep{gerganov2023llamacpp}. Two llama.cpp forks read
its 8-bit cache in place in the tensor-core FlashAttention kernel, with
Turing tiles, though neither reports a Turing measurement
\citep{piggidragon2026mmaq,twoxforever2026q8}. Its speedup over Triton is
expected of any tensor-core kernel on Turing, since Triton's runs on scalar
units there.

\textbf{Other optimizations and efficiency.} Running exllamav3's int8 path on 64\,KB Turing parts,
the host's choice of slice height and int8 verify rows have prior art
(Section~\ref{sec:build}), and the split-K kernel is rafatxf's. Comparable speedups over stock engines have been reported: rafatxf's branch reports
1.86 times at 4K and 2.57 times at 64K over exllamav3 1.5.3 on a 22\,GiB card at
4.0 bits per weight in its documentation \citep{rafatxf2026turing}, against
1.6 to 2.2 times in its upstream issue \citep{rafatxf2026issue417}, and a GTX~16 branch reports
1.93 times for another model \citep{nalimcos2026gtx16}; ours comes from
partly different mechanisms and a newer baseline. We make no efficiency claim. Our
absolute speed on the 11\,GiB card comes largely from a 2.0 bits-per-weight
quantization that fits it; reading a token's bytes at the card's streaming
rate would take \PeBytesErrMin--\PeBytesErrMax\% less time than the token
takes, while earlier Turing engines report 70--85\% of the bandwidth at 3.5 to
4 bits per weight on 22--24\,GiB cards
\citep{ninfer2026titan,rafatxf2026turing,vcruz3052026exl3turing}. Against
NInferno on this card the engines are close per byte. A token reads
\PeTokGbExl\,GB of EXL3 weights and \PeTokGbIxxs\,GB of UD-IQ2\_XXS (all
but the embedding table, which decoding reads one row at a time; the GGUF file
is \PeFileGbIxxs\,GB). In the 4~November confirmation session, w13's undrafted rate moved
about \PeCfGbsW\,GB/s and NInferno's about \PeCfGbsNif\,GB/s. Those dated rates imply similar weight-byte throughput; this arithmetic
is not a bandwidth measurement or an explanation of the replacement margin. UD-IQ2\_XXS reads
\PeTokBytesPct\% more weight bytes per token than EXL3, which accounts for
part of the dated undrafted rate difference.

\subsection{Configuration selection}
\label{sec:choice}
For short prompts, drafting with the MTP layer in mode~2 is the fastest
configuration we measured, and on GSM8K it answered every item as the stock
engine did. Up to \PeReachEE, the 8-bit cache read by fd75d keeps the decode
rate at depth (its accuracy was tested at 32K and 64K, not at 48K), and on
fresh items the split cache showed no accuracy advantage over it. The 8-bit cache failed at 80K (73--79K were not tried); past it
only the split and 4-bit caches fit: the split cache
reached \PeReachEF, where accuracy on this task fell, and the 4-bit cache was
not shown usable at 64K. For one question over a long prompt with a short
answer, prefill dominates: it ran at about \PeArmEESixtyFourPrefill{} tokens
per second at 64K, and a 64K request took \PeArmEESixtyFourItemS{} seconds on
average with fd75d's 8-bit cache. The decode rate matters when answers are
long or follow-up turns reuse a cached prompt. An fd75d path for the split
cache would carry the decode rate past \PeReachEE; it is future work.

\subsection{Detailed limitations and threats to validity}
\label{sec:limitations}
\textbf{Sessions.} Sessions of this harness differ by up to
\PeAuditSlowHi\% in time per token, so only comparisons within a session are
claims. F2 and the confirmation session agree within 1\%, but both ran on one
card in one desktop. The speed workload is three prompts of 256 sampled
tokens each, without a seed, and other programs raised the desktop's load
during F2. Rotation within a session spreads both across the arms, but the
absolute rates still reflect both.

\textbf{Competitors measured.} Each competitor ran the configuration
registered for it, or the first rung of the fit ladder that fit, with our
prompts, requested sampling knobs and timing. Engine-specific defaults and
sampling order remained different; neutral penalties in the replacement do
not establish matching distributions or quality. No competitor was tuned
further for these sessions. The replacement completed only
\PeFreshComplete{} of \PeFreshPlanned{} comparisons, with
\PeFreshEligible{} eligible current leads. Its incomplete stock cells,
unknown llama.cpp placement and unmeasured DFlash2 cells are not losses.
The following limitations describe the older confirmation sessions:
\begin{itemize}\setlength{\itemsep}{1pt}
\item rafatxf ran its 4-bit cache against our 8-bit one, and no same-cache arm
ran;
\item NInferno's DFlash2 lane ran on NInferno's own drafter-attached
conversion, at its automatic context when empty and on a survival rung at
28K, in the DFlash2 sessions rather than the confirmation session;
\item two drafted llama.cpp-family lanes did not fit at full offload;
\item llama.cpp's MTP lane is a partial-offload speed;
\item buun-llama-cpp's undrafted lane has one round;
\item formats that need a download, such as ik\_llama.cpp's trellis
quantizations, were not run.
\end{itemize}
The per-answer comparison with llama.cpp ran in other sessions, on builds of
September and October. Public reports we cite were measured elsewhere.

\textbf{Accuracy of the faster builds.} w13's accuracy rests on GSM8K with
MTP (identical outcomes to the stock wheel), on bit-identity for undrafted
decoding, and on K1 at 64K. Among the exllamav3 builds in the original 2080 Ti study, MMLU tested only the stock wheel. The cross-card replication in Appendix~\ref{app:fleet} tested w13 under different deployment conditions. The comparisons
with llama.cpp are descriptive; each lane ran once, and the MMLU gap to
UD-IQ2\_XXS is not significant after adjustment.

\textbf{One deployed configuration.} Apart from the laptop-GPU measurements
that Section~\ref{sec:bytes} cites, the results cover one card, one model at
one weight precision and single-stream requests, with sampled decoding for
speed and greedy decoding for accuracy and long context. Drafting
makes most decode steps run attention on several query tokens, which the
logit check, run one token at a time on gold answers, did not cover. fd75d's
accuracy is established at 64K (and, combined with the format change, at
32K); at 48K, where its logits diverged by the registered rule, no accuracy arm
ran. The 1-bit results of Section~\ref{sec:bytes} come from our patched
kernels on the 2080 Ti and the stock kernels on the laptop GPU, so the
difference between the two cards confounds the card with the kernels.

\textbf{Long-context design.} The graded task is one family, multi-value
retrieval with \PeNeedles{} values, scored all-or-nothing on chat-session
haystacks; it is a harder variant of a RULER task and not comparable with
RULER scores. fp16 was tested only to \PeReachFp, so every 64K comparison is against
the 8-bit cache, and at 96K only the split cache ran. Each comparison rests on
\PeHTwoN{} items, so where arms disagree on a fifth of the items the 90\%
interval is \PeWidthMin--\PeWidthMax{} points wide. The change of graded task
was registered before any test item ran, but after seeing that the planned
families could not separate the formats. Reach is the longest ladder step
that passed with one filled prompt and a short answer; for three formats it
is a lower bound.


\section{Registered predictions and their outcomes}
\label{sec:predictions}
Every prediction registered for this paper's measurements, with its outcome
(Tables~\ref{tab:predictions} and~\ref{tab:predictions-context}).
Failed predictions are kept: several of them changed what the paper claims.
The registrations, with their amendments and dated outcomes, ship in the
artifact, as does the separate registration of the kernel bench that
Section~\ref{sec:bytes} cites, with its own predictions and outcomes.

The laptop-GPU repeat that Section~\ref{sec:bytes} cites registered nine
predictions over three runs: at 80\,W, a replicate at 80\,W in the
performance power profile, and at 175\,W. Five held and four failed. As
predicted, the fixed parts were most of the calls' time above their reads in
all three runs, and the raised limit lifted the memory clock to
\PeAdaBoostMemMhz\,MHz. Against the predictions, at 80\,W a 1-bit call took
up to \PeAdaOneFasterHi\% less time than a 2-bit call in the same role,
against a predicted 10\% at most; the performance profile alone left decode's
memory clock at \PeAdaMemMhz\,MHz, not \PeAdaBoostMemMhz; and the faster
memory did not shrink the 1-bit advantage in either later run, which at
175\,W grew to \PeAdaBoostOneFasterLo--\PeAdaBoostOneFasterHi\% less time.

\begin{table}[!tbp]\centering\footnotesize\setlength{\tabcolsep}{4pt}
\begin{tabular}{@{}>{\raggedright\arraybackslash}p{3.75in}>{\raggedright\arraybackslash}p{1.95in}l@{}}
\toprule
Prediction & Outcome & Section\\\midrule
\multicolumn{3}{@{}l}{\textit{Confirmation sessions (4--5 November)}}\\
w13 at 43--48 tok/s undrafted and 64--72 with MTP at an empty context, in the registered session & held & \ref{sec:speed}\\
w13 over stock: 1.75--1.95$\times$ undrafted, 1.45--1.65$\times$ with MTP, $\geq$1.9$\times$ undrafted at 28K & held & \ref{sec:speed}\\
exllamav3 1.6.6 within $\pm5\%$ of 1.6.5 & held & \ref{app:engine-comparisons}\\
w13 over rafatxf's fork: $+3$ to $+12\%$ empty, $-3$ to $+10\%$ at 28K & held & \ref{app:engine-comparisons}\\
NInferno UD-IQ2\_XXS at 42--48 tok/s; w13 ahead by 0--12\% & held & \ref{app:engine-comparisons}\\
NInferno UD-IQ2\_S 8--18\% slower than UD-IQ2\_XXS & failed: 6.5\% slower & \ref{app:engine-comparisons}\\
llama.cpp with DFlash2 at 50--65 tok/s; with MTP at 40--55 & failed: no full-offload fit & \ref{app:engine-comparisons}\\
ik\_llama.cpp 0--25\% faster than llama.cpp & failed: 19.6\% slower & \ref{app:engine-comparisons}\\
The ``fastest'' sentence survives every session & held then; no current ranking & \ref{sec:new}\\
w13 with MTP leads NInferno with DFlash2 in every round, by 20--50\% empty and 5--35\% at 28K & held & \ref{app:engine-comparisons}\\
NInferno's DFlash2 fits the registered 32K context & failed & \ref{sec:speedsetup}\\
NInferno's DFlash2 cannot start at 29,696 tokens & failed: it started on a survival rung & \ref{sec:speedsetup}\\
\multicolumn{3}{@{}l}{\textit{Earlier release refresh (12 November; NInferno comparison replaced)}}\\
w13 over NInferno 65e573fd undrafted: 0.95--1.10$\times$ empty, 1.00--1.20$\times$ at 28K & held: \PeRfXNifNdZero, \PeRfXNifNdDeep & \ref{app:engine-comparisons}\\
w13 with MTP over NInferno's DFlash2: 1.05--1.35$\times$ empty, 1.10--1.45$\times$ at 28K & held: \PeRfXNifDfZero, \PeRfXNifDfDeep & \ref{app:engine-comparisons}\\
w13 at least 1.6$\times$ llama.cpp dd6b945b undrafted & held: \PeRfXLcppNdZero, \PeRfXLcppNdDeep; residency unrecorded & \ref{app:engine-comparisons}\\
\multicolumn{3}{@{}l}{\textit{Replacement comparison (14 November)}}\\
Four NInferno pooled leads hold in every round & held; neutral penalties, sampling order differs & \ref{app:engine-comparisons}\\
Every tested lead revalidated & not established: \PeFreshComplete/\PeFreshPlanned{} complete, \PeFreshEligible{} eligible & \ref{app:engine-comparisons}\\
\multicolumn{3}{@{}l}{\textit{Fleet deployment replication (13--14 November)}}\\
MTP faster at both fills on each card; 4080 SUPER faster than 4070 on the mixture; GSM8K MTP within 5 points of archived w13; within-card nd/MTP accuracy gaps $\leq3$ points & all \PeFleetAssessedPredictions{} held, descriptive & \ref{app:fleet}\\
\bottomrule\end{tabular}
\caption{Registered predictions and outcomes for the speed sessions and
cross-card deployment replication. Earlier and replacement comparisons retain
their dates and scope.}
\label{tab:predictions}
\end{table}

\begin{table}[!tbp]\centering\footnotesize\setlength{\tabcolsep}{4pt}
\begin{tabular}{@{}>{\raggedright\arraybackslash}p{3.75in}>{\raggedright\arraybackslash}p{1.95in}l@{}}
\toprule
Prediction & Outcome & Section\\\midrule
\multicolumn{3}{@{}l}{\textit{Long context}}\\
Peak memory within $-150$ to $+250$\,MiB of projection & failed for the quantized caches & \ref{sec:reach}\\
8-bit cache passes 64K; 4-bit cache reaches $\geq$128K & held & \ref{sec:reach}\\
H1 (fd75d's 8-bit cache against fp16, 32K) equivalent & held & \ref{sec:accuracy}\\
H2 (4-bit cache at 64K) usable and equivalent & failed: not shown usable & \ref{sec:accuracy}\\
H3 (4-bit cache at 32K) equivalent & failed: usable, not equivalent & \ref{sec:accuracy}\\
fd75d's logits do not diverge from Triton's at 28K, 48K and 64K & failed at 48K & \ref{sec:kernel}\\
K1: Triton reading the 8-bit cache is more accurate than fd75d at 64K & failed: equivalent & \ref{sec:kernel}\\
K2: 4-bit values cost nothing at 64K (equivalent) & failed: $+11.5$ & \ref{sec:rep}\\
R1: the split cache's 64K lead replicates on fresh items & failed: not replicated & \ref{sec:rep}\\
N: a harmless reordering changes at most 6 of 96 outcomes & held & \ref{sec:rep}\\
R2 and R3 equivalent & failed & \ref{sec:rep}\\
\multicolumn{3}{@{}l}{\textit{Mechanism}}\\
Bytes alone underpredict w13's undrafted F2 token at an empty context by 37--45\% & held & \ref{sec:bytes}\\
The priced graph within $\pm12\%$ of every held-out configuration & failed & \ref{sec:bytes}\\
The priced graph, extended to drafting, within $\pm10\%$ for w13 mode~2 & failed empty, held at 28K & \ref{sec:bytes}\\
Mode~2 ahead of rafatxf's fork by $\geq$10\% in F2 & failed: 8.2\% & \ref{app:engine-comparisons}\\
fd75d makes the 28K token $\geq$20\% faster, every round & held & \ref{sec:depth}\\
w11b within $\pm3\%$ of w9 at an empty context & failed: $+4.2\%$ & \ref{sec:depth}\\
W12p's products bit-identical for 1--7 rows & failed (narrowed to 1--2 rows, which held) & \ref{sec:build}\\
Accuracy checks of w11b and w13 mode~2 within $-1.0$ point & held & \ref{sec:quality}\\
\bottomrule\end{tabular}
\caption{Registered predictions and outcomes for long context and the
mechanism. The laptop-GPU repeat's nine predictions are given in the text.}
\label{tab:predictions-context}
\end{table}

\clearpage
\section{Prompts}
\label{app:prompts}
\textbf{Speed.} The speed harness decodes from three fixed prompts, each sent
as one user message through Qwen's chat template with thinking off (an empty
think block opens the assistant turn):
\begin{quote}\small
\textbf{Code.} Write a Python function that returns the n-th Fibonacci number
iteratively, with a docstring.\\
\textbf{Math.} A train travels 180 km in 2.5 hours, then 120 km in 1.5 hours.
What is its average speed? Show the steps.\\
\textbf{Prose.} Describe, in two paragraphs, how a lighthouse keeper's day
might have looked in 1890.
\end{quote}
The 28K arms put a fixed synthetic log ahead of the same prompt in the same
message. It opens with ``Here is a log for reference; the question after it
does not depend on it.'' and a blank line. Then come lines such as ``Record
00000: basalt nectar 78357 fjord batch 25.'', drawn from twenty fixed words by
a generator seeded with 1030. The text is cut to exactly 28,000 tokens of the
model's tokenizer and followed by a blank line. \q{tools/turing/census.py}
(\q{filler}) rebuilds it.
Table~\ref{tab:f2prompts} gives F2's rates by prompt.
% Generated by paper8/scripts/build_claims.py from testsuite/results/turing/ab/final2.jsonl. Do not edit by hand.
\begin{table}[!htbp]
\centering\small
\begin{tabular}{lllrrrrc}
\toprule
Build & Context & Drafting & Code & Math & Prose & Pooled & Tokens decoded\\
\midrule
stock & 0 & none & 25.4 & 25.3 & 25.3 & 25.3 & 699/768/768\\
w11b & 0 & none & 45.5 & 45.5 & 45.4 & 45.4 & 768/768/768\\
w13 & 0 & none & 47.6 & 47.6 & 47.5 & 47.5 & 768/768/768\\
w11b & 28K & none & 41.9 & 41.9 & 41.9 & 41.9 & 655/768/768\\
w13 & 28K & none & 43.9 & 43.9 & 43.9 & 43.9 & 551/765/768\\
stock & 0 & MTP & 59.5 & 59.5 & 32.9 & 46.4 & 654/753/756\\
w11b & 0 & MTP & 65.5 & 68.6 & 39.9 & 54.6 & 757/768/768\\
w13, mode 1 & 0 & MTP & 80.7 & 82.4 & 47.1 & 65.5 & 760/768/768\\
w13, mode 2 & 0 & MTP & 87.5 & 92.5 & 50.8 & 71.6 & 748/768/758\\
rafatxf & 0 & MTP & 83.1 & 82.3 & 48.3 & 66.2 & 678/744/768\\
w11b & 28K & MTP & 61.7 & 64.6 & 36.7 & 50.5 & 697/748/768\\
w13, mode 1 & 28K & MTP & 74.6 & 78.7 & 44.5 & 61.3 & 683/768/768\\
w13, mode 2 & 28K & MTP & 80.4 & 85.2 & 48.9 & 66.6 & 638/768/768\\
rafatxf & 28K & MTP & 77.7 & 79.8 & 45.9 & 63.1 & 675/768/768\\
\bottomrule
\end{tabular}
\caption{F2's decode rates by prompt: tokens per second, each prompt pooled over its three rounds, with the registered rate pooled over all three prompts, and the tokens each prompt decoded over its three runs (code/math/prose, of 768 each; a run ends early at the end-of-sequence token). The rafatxf arms ran on that fork's 4-bit cache. Exploratory: the registered comparisons pool the prompts.}
\label{tab:f2prompts}
\end{table}


\textbf{Long context.} Every request is one user message, sent with Qwen's
thinking mode off: the instruction below, a blank line, the haystack with its
needles, a blank line, and ``Question:'' followed by the family's question. The
instruction is the earlier study's, unchanged:
\begin{quote}\small
Below is a long text made of chat conversations. A few special facts are
hidden in it. Read the text, then answer the question that follows it. Reply
with the answer only.
\end{quote}
Each graded needle reads ``One of the special magic numbers for \slot{key}
is: \slot{value}.'', following RULER's wording \citep{hsieh2024ruler}, with
\PeNeedles{} values for one key, and the question is ``What are all the
special magic numbers for \slot{key} mentioned in the provided text?''. The
check set uses the earlier study's single-needle, multi-key and
variable-tracking questions with the same instruction.

\end{document}
